% Generated mechanically from the frozen smoothing.tex target.
% Do not edit this generated file; edit the bound locale source instead.

% OK, start here.
%

\setcounter{chapter}{15}
\chapter{环同态的光滑化}



\phantomsection
\label{smoothing-section-phantom}


\section{引言}
\label{smoothing-section-introduction}

\noindent
本章的主要结论如下：
$$
\fbox{诺特环之间的正则同态是光滑同态的滤过余极限。}
$$
该定理由 Popescu 证明，见 \cite{popescu-letter}。
Richard Swan 给出了 Popescu 证明的一份易读阐述，见 \cite{swan}；
他采用了 Andr\'e 的笔记以及 Ogoma 的一篇论文，见 \cite{Ogoma}。

\medskip\noindent
我们的叙述沿循 Swan 的处理，不过会先证明一个中间结果，
以便在稍为简单的情形中展开论证。下面概述全章结构。
首先解决如下“提升问题”：光滑代数的滤过余极限的平坦无穷小形变，
仍是光滑代数的滤过余极限。这个结果实质上说明，只须对约化诺特环之间的
同态证明主定理。随后证明两个十分巧妙的引理，分别称为“提升引理”和
“消奇异化引理”。我们将说明，合用这两个引理后，主定理可归结为证明：
域 $k$ 上几何正则的诺特代数 $\Lambda$ 是光滑 $k$-代数的滤过余极限。
接着讨论 Popescu--Ogoma--Swan--Andr\'e 证明所需的局部技巧。
最后三节给出证明。

\medskip\noindent
最后列出若干参考文献。设 $A$ 为 Hensel 诺特局部环。
若对任意 $f_1, \ldots, f_m \in A[x_1, \ldots, x_n]$，方程组
$f_1 = 0, \ldots, f_m = 0$ 在 $A$ 的完备化中有解当且仅当它在 $A$ 中有解，
则称 $A$ 具有{\it 逼近性质}。这一定义由 Artin 提出。
Artin 首先对解析方程组证明了逼近性质，见
\cite{Artin-Analytic-Approximation}。
在 \cite{Artin-Algebraic-Approximation} 中，Artin 对本质有限型于
优秀离散赋值环的局部环证明了逼近性质。
Artin 猜想（见 \cite{Artin-Algebraic-Approximation} 第 26 页），
每个优秀 Hensel 局部环都应具有逼近性质。

\medskip\noindent
后来人们提出猜想：诺特环的每个正则同态都是光滑代数的滤过余极限（例如见
\cite{Raynaud-Rennes}, \cite{popescu-global}, \cite{Artin-power-series},
\cite{Artin-Denef}）。我们不确定这一猜想应归于谁。\footnote{在
\cite{Artin-power-series}、\cite{Artin-Denef} 和 \cite{popescu-global}
中表述的问题或猜想更强；\cite{Cipu} 证明了它与原版本等价。}
它与逼近性质的联系如下：若 $A$ 如上，并且 $A \to A^\wedge$ 是光滑代数的
余极限，则逼近性质成立（此处待补引用）。此外，若 $A$ 是优秀局部环，
则主定理适用于同态 $A \to A^\wedge$，因为优秀局部环的条件之一是其
形式纤维几何正则。优秀局部环由 Grothendieck 定义，该定义于 1965 年刊印。

\medskip\noindent
\cite{Artin-power-series} 证明了：对任意本质有限型于域的局部环 $R$，
$R \to R^\wedge$ 都是光滑代数的滤过余极限。
\cite{Rotthaus-Artin} 证明了：对任意本质有限型于优秀离散赋值环的局部环
$R$，$R \to R^\wedge$ 都是光滑代数的滤过余极限。
最后，如上所述，主定理在
\cite{popescu-GND}, \cite{popescu-GNDA}, \cite{popescu-letter}, 
\cite{Ogoma} 和 \cite{swan} 中得到证明。

\medskip\noindent
反过来，利用上述若干结果，\cite{Rotthaus-excellent} 证明了：
任何具有逼近性质的诺特局部环都是优秀环。

\medskip\noindent
论文 \cite{Spivakovsky} 给出了主定理的另一种处理方式，但似乎颇难读懂
（例如，\cite[Lemma 5.2]{Spivakovsky} 看来是对 \cite[Lemma 3]{Elkik}
的不正确改述）。关于这些材料还有一篇 Bourbaki 报告，见 \cite{Teissier}。












\section{奇异理想}
\label{smoothing-section-singular-ideal}

\noindent
设 $R \to A$ 为环同态。$A$ 相对于 $R$ 的奇异理想，是 $A$ 中定义态射
$\Spec(A) \to \Spec(R)$ 的奇异轨迹的根理想。其严格定义如下。

\begin{definition}
\label{smoothing-definition-singular-ideal}
设 $R \to A$ 为环同态。$A$ 相对于 $R$ 的{\it 奇异理想}记作 $H_{A/R}$，
它是满足下式的唯一根理想 $H_{A/R} \subset A$：
$$
V(H_{A/R}) = \{\mathfrak q \in \Spec(A) \mid R \to A
\text{ 在 }\mathfrak q\text{ 处不光滑}\}.
$$
\end{definition}

\noindent
这一定义合理，因为 $R \to A$ 光滑的素理想所成集合是开集，见
《代数》定义 \ref{algebra-definition-smooth-at-prime}。
为了找出奇异理想的一组显式生成元，我们先证明以下引理。

\begin{lemma}
\label{smoothing-lemma-find-strictly-standard}
设 $R$ 为环，$A = R[x_1, \ldots, x_n]/(f_1, \ldots, f_m)$，
$\mathfrak q \subset A$ 为素理想。假设 $R \to A$ 在 $\mathfrak q$ 处光滑。
则存在 $a \in A$ 且 $a \not \in \mathfrak q$、整数 $c$ 且
$0 \leq c \leq \min(n, m)$，以及基数均为 $c$ 的子集
$U \subset \{1, \ldots, n\}$、$V \subset \{1, \ldots, m\}$，使得
$$
a = a' \det(\partial f_j/\partial x_i)_{j \in V, i \in U}
$$
对某个 $a' \in A$ 成立，并且
$$
a f_\ell \in (f_j, j \in V) + (f_1, \ldots, f_m)^2
$$
对所有 $\ell \in \{1, \ldots, m\}$ 都成立。
\end{lemma}

\begin{proof}
令 $I = (f_1, \ldots, f_m)$。于是 $A$ 在 $R$ 上的朴素余切复形同伦等价于
$I/I^2 \to \bigoplus A\text{d}x_i$，见
《代数》引理 \ref{algebra-lemma-NL-homotopy}。
我们将用到朴素余切复形的构造与局部化可交换；见
《代数》第 \ref{algebra-section-netherlander} 节，特别是
《代数》引理 \ref{algebra-lemma-localize-NL}。
由《代数》定义 \ref{algebra-definition-smooth} 和
\ref{algebra-definition-smooth-at-prime} 可知，对某个
$a \in A$、$a \not \in \mathfrak q$，同态
$(I/I^2)_a \to \bigoplus A_a\text{d}x_i$ 是分裂单射。
重排 $x_1, \ldots, x_n$ 与 $f_1, \ldots, f_m$ 的编号后，可以假设
$f_1, \ldots, f_c$ 构成向量空间
$I/I^2 \otimes_A \kappa(\mathfrak q)$ 的一组基，并且
$\text{d}x_{c + 1}, \ldots, \text{d}x_n$ 的像构成
$\Omega_{A/R} \otimes_A \kappa(\mathfrak q)$ 的一组基。
因此，用 $aa'$ 代替 $a$（其中 $a' \in A$ 且
$a' \not \in \mathfrak q$）后，可以假设 $f_1, \ldots, f_c$
构成 $(I/I^2)_a$ 的一组基，并且
$\text{d}x_{c + 1}, \ldots, \text{d}x_n$ 的像构成
$(\Omega_{A/R})_a$ 的一组基。在此情形中，当整数 $N$ 充分大时，
$a^N$ 满足引理的条件（取 $U = V = \{1, \ldots, c\}$）。
\end{proof}

\noindent
我们将使用 $A$ 中相对于 $R$ 的{\it 严格标准}元这一概念。
这里的概念比 Swan 的论文 \cite{swan} 中所用的概念稍弱。
我们还把上一个引理中所得类型的元素定义为{\it 初等标准}元。
引理 \ref{smoothing-lemma-compare-standard} 将比较这些不同类型的元素。

\begin{definition}
\label{smoothing-definition-strictly-standard}
设 $R \to A$ 为有限表示的环同态，$a \in A$。若存在一个表示
$A = R[x_1, \ldots, x_n]/(f_1, \ldots, f_m)$
以及满足 $0 \leq c \leq \min(n, m)$ 的整数 $c$，使得
\begin{equation}
\label{smoothing-equation-elementary-standard-one}
a = a' \det(\partial f_j/\partial x_i)_{i, j = 1, \ldots, c}
\end{equation}
对某个 $a' \in A$ 成立，并且
\begin{equation}
\label{smoothing-equation-elementary-standard-two}
a f_{c + j} \in (f_1, \ldots, f_c) + (f_1, \ldots, f_m)^2
\end{equation}
对 $j = 1, \ldots, m - c$ 成立，则称 $a$ 是
{\it $A$ 中相对于 $R$ 的初等标准元}。若存在一个表示
$A = R[x_1, \ldots, x_n]/(f_1, \ldots, f_m)$
以及满足 $0 \leq c \leq \min(n, m)$ 的整数 $c$，使得
\begin{equation}
\label{smoothing-equation-strictly-standard-one}
a = \sum\nolimits_{I \subset \{1, \ldots, n\},\ |I| = c}
a_I \det(\partial f_j/\partial x_i)_{j = 1, \ldots, c,\ i \in I}
\end{equation}
对某些 $a_I \in A$ 成立，并且
\begin{equation}
\label{smoothing-equation-strictly-standard-two}
a f_{c + j} \in (f_1, \ldots, f_c) + (f_1, \ldots, f_m)^2
\end{equation}
对 $j = 1, \ldots, m - c$ 成立，则称 $a$ 是
{\it $A$ 中相对于 $R$ 的严格标准元}。
\end{definition}

\noindent
以下引理可用来理解式 (\ref{smoothing-equation-strictly-standard-one}) 所蕴含的结论。

\begin{lemma}
\label{smoothing-lemma-parse-equation-strictly-standard-one}
设 $R$ 为环，$A = R[x_1, \ldots, x_n]/(f_1, \ldots, f_m)$，并记
$I = (f_1, \ldots, f_m)$。设 $a \in A$。则式
(\ref{smoothing-equation-strictly-standard-one}) 蕴含存在一个 $A$-线性同态
$\psi : \bigoplus\nolimits_{i = 1, \ldots, n} A \text{d}x_i \to A^{\oplus c}$
使得复合
$$
A^{\oplus c} \xrightarrow{(f_1, \ldots, f_c)}
I/I^2 \xrightarrow{f \mapsto \text{d}f}
\bigoplus\nolimits_{i = 1, \ldots, n} A \text{d}x_i
\xrightarrow{\psi}
A^{\oplus c}
$$
就是乘以 $a$ 的同态。反之，若这样的 $\psi$ 存在，则 $a^c$ 满足
(\ref{smoothing-equation-strictly-standard-one})。
\end{lemma}

\begin{proof}
这是《代数》引理 \ref{algebra-lemma-matrix-left-inverse} 的一个特例。
\end{proof}

\begin{lemma}[Elkik]
\label{smoothing-lemma-elkik}
设 $R \to A$ 为有限表示的环同态。奇异理想 $H_{A/R}$ 既是由
$A$ 中相对于 $R$ 的严格标准元生成的理想之根，也是由
$A$ 中相对于 $R$ 的初等标准元生成的理想之根。
\end{lemma}

\begin{proof}
假设 $a$ 是 $A$ 中相对于 $R$ 的严格标准元。我们断言 $A_a$ 在 $R$ 上光滑，
这便证明了 $a \in H_{A/R}$。事实上，令
$A = R[x_1, \ldots, x_n]/(f_1, \ldots, f_m)$、$c$ 和 $a' \in A$
如定义 \ref{smoothing-definition-strictly-standard} 所述。记
$I = (f_1, \ldots, f_m)$，则 $A$ 在 $R$ 上的朴素余切复形为
$I/I^2 \to \bigoplus A\text{d}x_i$。
条件 (\ref{smoothing-equation-strictly-standard-two}) 蕴含 $(I/I^2)_a$
由 $f_1, \ldots, f_c$ 的类生成。条件
(\ref{smoothing-equation-strictly-standard-one}) 蕴含微分同态
$(I/I^2)_a \to \bigoplus A_a\text{d}x_i$ 有左逆，见
引理 \ref{smoothing-lemma-parse-equation-strictly-standard-one}。
因此，由定义和《代数》引理 \ref{algebra-lemma-localize-NL}，
$R \to A_a$ 光滑。

\medskip\noindent
令 $H_e, H_s \subset A$ 分别为由 $A$ 中相对于 $R$ 的初等标准元、
严格标准元生成的理想之根。由定义和刚才的证明，
$H_e \subset H_s \subset H_{A/R}$。包含关系 $H_{A/R} \subset H_e$
由引理 \ref{smoothing-lemma-find-strictly-standard} 得出。
\end{proof}

\begin{example}
\label{smoothing-example-not-quasi-compact}
有限表示环同态的光滑点集未必是拟紧开集。例如，令
$R = k[x, y_1, y_2, y_3, \ldots]/(xy_i)$ 且 $A = R/(x)$。
则 $R \to A$ 的光滑轨迹为 $\bigcup D(y_i)$，它不是拟紧的。
\end{example}

\begin{lemma}
\label{smoothing-lemma-strictly-standard-base-change}
设 $R \to A$ 为有限表示的环同态，$R \to R'$ 为环同态。
若 $a \in A$ 是 $A$ 中相对于 $R$ 的初等标准元（相应地，严格标准元），
则 $a \otimes 1$ 是 $A \otimes_R R'$ 中相对于 $R'$ 的初等标准元
（相应地，严格标准元）。
\end{lemma}

\begin{proof}
若 $A = R[x_1, \ldots, x_n]/(f_1, \ldots, f_m)$ 是 $A$ 在 $R$ 上的一个表示，
则
$A \otimes_R R' = R'[x_1, \ldots, x_n]/(f'_1, \ldots, f'_m)$
是 $A \otimes_R R'$ 在 $R'$ 上的一个表示，其中 $f'_j$ 是 $f_j$
在 $R'[x_1, \ldots, x_n]$ 中的像。因此结论由定义立即得出。
\end{proof}

\begin{lemma}
\label{smoothing-lemma-final-solve}
设 $R \to A \to \Lambda$ 为环同态，且 $A$ 是有限表示 $R$-代数。
假设 $H_{A/R} \Lambda = \Lambda$。则存在分解
$A \to B \to \Lambda$，其中 $B$ 在 $R$ 上光滑。
\end{lemma}

\begin{proof}
选取 $f_1, \ldots, f_r \in H_{A/R}$ 和
$\lambda_1, \ldots, \lambda_r \in \Lambda$，使得在 $\Lambda$ 中
$\sum f_i\lambda_i = 1$。令
$B = A[x_1, \ldots, x_r]/(f_1x_1 + \ldots + f_rx_r - 1)$
并以 $x_i \mapsto \lambda_i$ 定义 $B \to \Lambda$。
为验证 $B$ 在 $R$ 上光滑，注意由 $H_{A/R}$ 的定义，$A_{f_i}$
在 $R$ 上光滑，而 $B_{f_i}$ 在 $A_{f_i}$ 上光滑。细节从略。
\end{proof}





\section{代数的表示}
\label{smoothing-section-presentations}

\noindent
本节的若干结果由 Elkik 证明。注意，下一个引理中的代数 $C$ 是 $A$ 上的
对称代数。此外，若 $R$ 是诺特环，则 $C$ 是有限表示 $R$-代数。

\begin{lemma}
\label{smoothing-lemma-improve-presentation}
设 $R$ 为环，$A$ 为有限表示 $R$-代数。存在具有缩回的有限型
$R$-代数同态 $A \to C$，且满足以下两个性质：
\begin{enumerate}
\item 对每个使 $R \to A_a$ 为局部完全交的 $a \in A$
（见《代数详论》定义
\ref{more-algebra-definition-local-complete-intersection}），
$C_a$ 在 $A_a$ 上光滑，并且有一个表示
$C_a = R[y_1, \ldots, y_m]/J$，其中 $J/J^2$ 是自由 $C_a$-模；
\item 对每个使 $A_a$ 在 $R$ 上光滑的 $a \in A$，模
$\Omega_{C_a/R}$ 是自由 $C_a$-模。
\end{enumerate}
\end{lemma}

\begin{proof}
选取一个表示 $A = R[x_1, \ldots, x_n]/I$，并记
$I = (f_1, \ldots, f_m)$。用短正合列
$$
0 \to K \to A^{\oplus m} \to I/I^2 \to 0
$$
定义 $A$-模 $K$；其中中间项的第 $j$ 个基向量 $e_j$ 映到右端 $f_j$ 的类。
令
$$
C = \text{Sym}^*_A(I/I^2).
$$
缩回就是投影到 $C$ 的 $0$ 次部分。存在满射
$R[x_1, \ldots, x_n, y_1, \ldots, y_m] \to C$，把 $y_j$ 映到
$I/I^2$ 中 $f_j$ 的类。此同态的核 $J$ 由 $f_1, \ldots, f_m$，
以及所有形如 $\sum h_j y_j$ 的元素生成；这里
$h_j \in R[x_1, \ldots, x_n]$，且 $\sum h_j e_j$ 定义 $K$ 的一个元素。
把《代数》引理 \ref{algebra-lemma-exact-sequence-NL} 应用于
$R \to A \to C$ 及上述表示，并使用《代数详论》引理
\ref{more-algebra-lemma-cotangent-complex-symmetric-algebra}，得到短正合列
\begin{equation}
\label{smoothing-equation-sequence}
I/I^2 \otimes_A C \to J/J^2 \to K \otimes_A C \to 0
\end{equation}
的 $C$-模。设 $h \in R[x_1, \ldots, x_n]$ 的像为 $a \in A$。
对局部化环采用如下表示：
$$
A_a = R[x_0, x_1, \ldots, x_n]/I'
\quad\text{且}\quad
C_a = R[x_0, x_1, \ldots, x_n, y_1, \ldots, y_m]/J'
$$
其中 $I' = (hx_0 - 1, I)$，$J' = (hx_0 - 1, J)$。因此，作为
$A_a$-模有 $I'/(I')^2 = A_a \oplus (I/I^2)_a$，作为 $C_a$-模有
$J'/(J')^2 = C_a \oplus (J/J^2)_a$。于是得到
\begin{equation}
\label{smoothing-equation-sequence-localized}
C_a \oplus I/I^2 \otimes_A C_a \to
C_a \oplus (J/J^2)_a \to
K \otimes_A C_a \to 0
\end{equation}
这正是《代数》引理 \ref{algebra-lemma-exact-sequence-NL} 对应于
$R \to A_a \to C_a$ 及上述表示的正合列。

\medskip\noindent
接下来，假设 $a \in A$ 使得 $A_a$ 是 $R$ 上的局部完全交。
则 $(I/I^2)_a$ 是有限投射 $A_a$-模，见《代数详论》引理
\ref{more-algebra-lemma-quasi-regular-ideal-finite-projective}。
因此 $K_a \oplus (I/I^2)_a \cong A_a^{\oplus m}$ 是自由模；特别地，
$K_a$ 也是有限投射模。由《代数详论》引理
\ref{more-algebra-lemma-transitive-lci-at-end}，正合列
(\ref{smoothing-equation-sequence-localized}) 在左端也正合。因此
$$
J'/(J')^2 \cong
C_a \oplus I/I^2 \otimes_A C_a \oplus K \otimes_A C_a \cong
C_a^{\oplus m + 1}
$$
这就证明了 (1)。最后，再假设 $A_a$ 在 $R$ 上光滑。同一表示表明，
$\Omega_{C_a/R}$ 是同态
$$
J'/(J')^2 \longrightarrow
\bigoplus\nolimits_i C_a\text{d}x_i \oplus \bigoplus\nolimits_j C_a\text{d}y_j
$$
的余核。在上述分解中，$J'/(J')^2$ 的直和项 $C_a$ 对应于
$hx_0 - 1$，因而同构地映到直和项 $C_a\text{d}x_0$。
$J'/(J')^2$ 的直和项 $I/I^2 \otimes_A C_a$ 单射到
$\bigoplus_{i = 1, \ldots, n} C_a\text{d}x_i$，其商为
$\Omega_{A_a/R} \otimes_{A_a} C_a$。直和项 $K \otimes_A C_a$
单射到 $\bigoplus_{j \geq 1} C_a\text{d}y_j$，其商同构于
$I/I^2 \otimes_A C_a$。所以，上一个显示同态的余核是模
$I/I^2 \otimes_A C_a \oplus \Omega_{A_a/R} \otimes_{A_a} C_a$。
由光滑环同态的定义，$(I/I^2)_a \oplus \Omega_{A_a/R}$ 是自由模，
从而 (2) 成立。
\end{proof}

\noindent
对 Hensel 对上的光滑环同态，以下命题由 Elkik 在 \cite{Elkik} 中证明。
关于光滑环同态的表述可见 \cite{Arabia}；该文还证明了光滑代数之间的
环同态可以提升。

\begin{proposition}
\label{smoothing-proposition-lift-smooth}
\begin{slogan}
光滑代数和 syntomic 代数可沿满射提升
\end{slogan}
设 $R \to R_0$ 是核为 $I$ 的满环同态。
\begin{enumerate}
\item 若 $R_0 \to A_0$ 是 syntomic 环同态，则存在 syntomic 环同态
$R \to A$，使得 $A/IA \cong A_0$。
\item 若 $R_0 \to A_0$ 是光滑环同态，则存在光滑环同态
$R \to A$，使得 $A/IA \cong A_0$。
\end{enumerate}
\end{proposition}

\begin{proof}
假设 $R_0 \to A_0$ 是 syntomic 的；特别地，它是局部完全交
（见《代数详论》引理 \ref{more-algebra-lemma-syntomic-lci}）。
选取表示 $A_0 = R_0[x_1, \ldots, x_n]/J_0$，并令
$C_0 = \text{Sym}^*_{A_0}(J_0/J_0^2)$。注意，$J_0/J_0^2$ 是有限投射
$A_0$-模（见《代数》引理
\ref{algebra-lemma-syntomic-presentation-ideal-mod-squares}）。
由引理 \ref{smoothing-lemma-improve-presentation}，环同态 $A_0 \to C_0$ 光滑，
并且可找到表示 $C_0 = R_0[y_1, \ldots, y_m]/K_0$，使得
$K_0/K_0^2$ 是自由 $C_0$-模。由《代数》引理
\ref{algebra-lemma-huber}，可以假设
$C_0 = R_0[y_1, \ldots, y_m]/(\overline{f}_1, \ldots, \overline{f}_c)$
其中 $\overline{f}_1, \ldots, \overline{f}_c$ 的像构成 $K_0/K_0^2$
在 $C_0$ 上的一组基。选取 $\overline{f}_1, \ldots, \overline{f}_c$
的提升 $f_1, \ldots, f_c \in R[y_1, \ldots, y_c]$，并令
$$
C = R[y_1, \ldots, y_m]/(f_1, \ldots, f_c)
$$
由构造，$C_0 = C/IC$。由《代数》引理
\ref{algebra-lemma-localize-relative-complete-intersection}，
用 $C_g$ 代替 $C$ 后，可以假设 $C$ 是 $R$ 上的相对整体完全交。
因此存在有限投射 $A_0$-模 $P_0$，使
$C_0 = \text{Sym}^*_{A_0}(P_0)$ 同构于某个 syntomic $R$-代数
$C$ 的 $C/IC$。

\medskip\noindent
选取整数 $n$ 以及直和分解
$A_0^{\oplus n} = P_0 \oplus Q_0$.
由《代数详论》引理 \ref{more-algebra-lemma-lift-projective-module}，
可找到平展环同态 $C \to C'$，它诱导同构 $C/IC \to C'/IC'$；
还可找到有限投射 $C'$-模 $Q$，使 $Q/IQ$ 同构于
$Q_0 \otimes_{A_0} C/IC$。于是 $D = \text{Sym}_{C'}^*(Q)$ 是光滑
$C'$-代数（见《代数详论》引理
\ref{more-algebra-lemma-symmetric-algebra-smooth}）。图示如下：
$$
\xymatrix{
R \ar[d] \ar[rr] & &
C \ar[r] \ar[d] &
C' \ar[r] \ar[d] &
D \ar[d] \\
R/I \ar[r] &
A_0 \ar[r] &
C/IC \ar[r]^{\cong} &
C'/IC' \ar[r] &
D/ID
}
$$
注意，$Q$ 的选取给出
\begin{align*}
D/ID & =
\text{Sym}_{C/IC}^*(Q_0 \otimes_{A_0} C/IC) \\
& =
\text{Sym}_{A_0}^*(Q_0) \otimes_{A_0} C/IC \\
& =
\text{Sym}_{A_0}^*(Q_0) \otimes_{A_0}
\text{Sym}_{A_0}^*(P_0) \\
& =
\text{Sym}_{A_0}^*(Q_0 \oplus P_0) \\
& =
\text{Sym}_{A_0}^*(A_0^{\oplus n}) \\
& =
A_0[x_1, \ldots, x_n]
\end{align*}
选取 $f_1, \ldots, f_n \in D$，使它们在
$D/ID = A_0[x_1, \ldots, x_n]$ 中分别映到 $x_1, \ldots, x_n$。
令 $A = D/(f_1, \ldots, f_n)$。注意 $A_0 = A/IA$。
我们断言，$R \to A$ 在 $V(IA)$ 的某个邻域上是 syntomic 的。
若断言成立，就可找到映到 $1 \in A_0$ 的 $f \in A$，使 $A_f$
在 $R$ 上 syntomic；于是 (1) 的证明完成。

\medskip\noindent
证明上述断言。注意，$R \to D$ 是 syntomic 环同态 $R \to C$、
平展环同态 $C \to C'$ 与光滑环同态 $C' \to D$ 的复合，因而是
syntomic 的（见《代数》引理
\ref{algebra-lemma-composition-syntomic} 以及
\ref{algebra-lemma-smooth-syntomic}）。这个问题在 $\Spec(D)$ 上是局部的，
故可假设 $D$ 是相对整体完全交（见《代数》引理
\ref{algebra-lemma-syntomic}）。设
$D = R[y_1, \ldots, y_m]/(g_1, \ldots, g_s)$。
令 $f'_1, \ldots, f'_n \in R[y_1, \ldots, y_m]$ 为
$f_1, \ldots, f_n$ 的提升。应用《代数》引理
\ref{algebra-lemma-localize-relative-complete-intersection} 即得断言。

\medskip\noindent
证明 (2)。由于光滑环同态是 syntomic 的，可找到 syntomic 环同态
$R \to A$，使得 $A_0 = A/IA$。由假设，$R \to A$ 的纤维在
$V(I)$ 中的素理想上光滑，故 $R \to A$ 在 $V(IA)$ 的一个开邻域上光滑
（见《代数》引理 \ref{algebra-lemma-flat-fibre-smooth}）。
因此，以一个局部化代替 $A$，便得到所需结论。
\end{proof}

\noindent
我们知道，任意 syntomic 环同态 $R \to A$ 在局部都是相对整体完全交，见
《代数》引理 \ref{algebra-lemma-syntomic}。下一个引理说明，
$\Spec(A)$ 上的向量丛是相对整体完全交。

\begin{lemma}
\label{smoothing-lemma-syntomic-complete-intersection}
设 $R \to A$ 为 syntomic 环同态。则存在具有缩回的光滑 $R$-代数同态
$A \to C$，使得 $C$ 是 $R$ 上的相对整体完全交，即
$$
C \cong R[x_1, \ldots, x_n]/(f_1, \ldots, f_c)
$$
在 $R$ 上平坦，且所有纤维的维数均为 $n - c$。
\end{lemma}

\begin{proof}
应用引理 \ref{smoothing-lemma-improve-presentation} 得到 $A \to C$。
由《代数》引理 \ref{algebra-lemma-huber}，可写成
$C = R[x_1, \ldots, x_n]/(f_1, \ldots, f_c)$，其中 $f_i$ 的像构成
$J/J^2$ 的一组基。环同态 $R \to C$ 是一个 syntomic 环同态与一个光滑
环同态的复合，因而是 syntomic 的（所以是平坦的）。由《代数》引理
\ref{algebra-lemma-lci}，纤维的维数为 $n - c$
（这些纤维是局部完全交，故该引理适用）。
\end{proof}

\begin{lemma}
\label{smoothing-lemma-smooth-standard-smooth}
设 $R \to A$ 为光滑环同态。则存在具有缩回的光滑 $R$-代数同态
$A \to B$，使得 $B$ 在 $R$ 上标准光滑，即
$$
B \cong R[x_1, \ldots, x_n]/(f_1, \ldots, f_c)
$$
并且 $\det(\partial f_j/\partial x_i)_{i, j = 1, \ldots, c}$
在 $B$ 中可逆。
\end{lemma}

\begin{proof}
应用引理 \ref{smoothing-lemma-syntomic-complete-intersection}，得到具有缩回的光滑
$R$-代数同态 $A \to C$，使得
$C = R[x_1, \ldots, x_n]/(f_1, \ldots, f_c)$ 是 $R$ 上的相对整体完全交。
由于 $C$ 在 $R$ 上光滑，存在短正合列
$$
0 \to
\bigoplus\nolimits_{j = 1, \ldots, c} C f_j \to
\bigoplus\nolimits_{i = 1, \ldots, n} C\text{d}x_i \to
\Omega_{C/R} \to 0
$$
由于 $\Omega_{C/R}$ 是投射 $C$-模，此正合列分裂。选取第一个同态的左逆
$t$。记
$t(\text{d}x_i) = \sum c_{ij} f_j$
于是 $\sum_i \frac{\partial f_j}{\partial x_i} c_{i\ell} = \delta_{j\ell}$
（Kronecker delta）。令
$$
B' = C[y_1, \ldots, y_c] =
R[x_1, \ldots, x_n, y_1, \ldots, y_c]/(f_1, \ldots, f_c)
$$
$R$-代数同态 $C \to B'$ 有一个由 $y_j \mapsto 0$ 给出的缩回。
我们断言，同态
$$
R[z_1, \ldots, z_n] \longrightarrow B',\quad
z_i \longmapsto x_i - \sum\nolimits_j c_{ij} y_j
$$
在 $\Spec(C) \to \Spec(B')$ 的像中每一点处都是平展的。
在 $\Omega_{B'/R[z_1, \ldots, z_n]}$ 中有
$$
0 =
\text{d}f_j - \sum\nolimits_i \frac{\partial f_j}{\partial x_i} \text{d}z_i
\equiv
\sum\nolimits_{i, \ell}
\frac{\partial f_j}{\partial x_i} c_{i\ell} \text{d}y_\ell
\equiv
\text{d}y_j \bmod (y_1, \ldots, y_c)\Omega_{B'/R[z_1, \ldots, z_n]}
$$
又因为模去
$\sum B'\text{d}y_j + (y_1, \ldots, y_c)\Omega_{B'/R[z_1, \ldots, z_n]}$
后有 $0 = \text{d}z_i = \text{d}x_i$，故
$$
\Omega_{B'/R[z_1, \ldots, z_n]}/
(y_1, \ldots, y_c)\Omega_{B'/R[z_1, \ldots, z_n]} = 0.
$$
由于 $\Omega_{B'/R[z_1, \ldots, z_n]}$ 是有限 $B'$-模，Nakayama 引理
给出某个 $g \in 1 + (y_1, \ldots, y_c)$，使得
$(\Omega_{B'/R[z_1, \ldots, z_n]})_g = 0$。这证明了
$R[z_1, \ldots, z_n] \to B'_g$ 非分歧，见《代数》定义
\ref{algebra-definition-unramified}。对任意以域 $k$ 为目标的环同态
$R \to k$，得到两个维数为 $n$ 的光滑 $k$-代数之间的非分歧环同态
$k[z_1, \ldots, z_n] \to (B'_g) \otimes_R k$。由《代数》引理
\ref{algebra-lemma-CM-over-regular-flat} 和
\ref{algebra-lemma-characterize-smooth-kbar}，
$k[z_1, \ldots, z_n] \to (B'_g) \otimes_R k$ 是平坦的。
再由逐纤维平坦性判别法（crit\`ere de platitude par fibre；见《代数》引理
\ref{algebra-lemma-criterion-flatness-fibre}），可知
$R[z_1, \ldots, z_n] \to B'_g$ 平坦。最后，《代数》引理
\ref{algebra-lemma-characterize-etale} 蕴含
$R[z_1, \ldots, z_n] \to B'_g$ 平展。令 $B = B'_g$。
注意 $C \to B$ 光滑且有缩回，故 $A \to B$ 也光滑且有缩回。
此外，$R[z_1, \ldots, z_n] \to B$ 平展。由《代数》引理
\ref{algebra-lemma-etale-standard-smooth}，可写成
$$
B = R[z_1, \ldots, z_n, w_1, \ldots, w_m]/(g_1, \ldots, g_m)
$$
其中 $\det(\partial g_j/\partial w_i)$ 在 $B$ 中可逆。引理得证。
\end{proof}

\begin{lemma}
\label{smoothing-lemma-colimit-standard-smooth}
设 $R \to \Lambda$ 为环同态。若 $\Lambda$ 是光滑 $R$-代数的滤过余极限，
则 $\Lambda$ 是标准光滑 $R$-代数的滤过余极限。
\end{lemma}

\begin{proof}
设 $A \to \Lambda$ 为 $R$-代数同态，其中 $A$ 是有限表示 $R$-代数。
按照《代数》引理 \ref{algebra-lemma-when-colimit}，只需把此同态经由某个
标准光滑代数分解。我们已知它可以分解为 $A \to B \to \Lambda$，其中
$B$ 在 $R$ 上光滑。由引理 \ref{smoothing-lemma-smooth-standard-smooth}，选取具有缩回
$C \to B$ 的 $R$-代数同态 $B \to C$，使 $C$ 在 $R$ 上标准光滑。
所求分解便是 $A \to B \to C \to B \to \Lambda$。
\end{proof}

\begin{lemma}
\label{smoothing-lemma-standard-smooth-include-generators}
设 $R \to A$ 为标准光滑环同态，$E \subset A$ 为有限子集，且 $|E| = n$。
则存在表示
$A = R[x_1, \ldots, x_{n + m}]/(f_1, \ldots, f_c)$，其中 $c \geq n$，
使 $\det(\partial f_j/\partial x_i)_{i, j = 1, \ldots, c}$ 在 $A$ 中可逆，
并且 $E$ 恰为 $x_1, \ldots, x_n$ 的同余类所成集合。
\end{lemma}

\begin{proof}
选取表示 $A = R[y_1, \ldots, y_m]/(g_1, \ldots, g_d)$，使得
$\det(\partial g_j/\partial y_i)_{i, j = 1, \ldots, d}$ 的像在 $A$ 中可逆。
选定枚举 $E = \{a_1, \ldots, a_n\}$，并选取
$h_i \in R[y_1, \ldots, y_m]$，使其在 $A$ 中的像为 $a_i$。
考虑表示
$$
A = R[x_1, \ldots, x_n, y_1, \ldots, y_m]/
(x_1 - h_1, \ldots, x_n - h_n, g_1, \ldots, g_d)
$$
并令 $c = n + d$。
\end{proof}

\begin{lemma}
\label{smoothing-lemma-compare-standard}
设 $R \to A$ 为有限表示环同态，$a \in A$。考虑关于 $a$ 的下列条件：
\begin{enumerate}
\item $A_a$ 在 $R$ 上光滑；
\item $A_a$ 在 $R$ 上光滑，且 $\Omega_{A_a/R}$ 稳定自由；
\item $A_a$ 在 $R$ 上光滑，且 $\Omega_{A_a/R}$ 自由；
\item $A_a$ 在 $R$ 上标准光滑；
\item $a$ 是 $A$ 中相对于 $R$ 的严格标准元；
\item $a$ 是 $A$ 中相对于 $R$ 的初等标准元。
\end{enumerate}
则有
\begin{enumerate}
\item[(a)] (4) $\Rightarrow$ (3) $\Rightarrow$ (2) $\Rightarrow$ (1),
\item[(b)] (6) $\Rightarrow$ (5),
\item[(c)] (6) $\Rightarrow$ (4),
\item[(d)] (5) $\Rightarrow$ (2),
\item[(e)] (2) $\Rightarrow$ 当 $e \geq e_0$ 时，$a^e$ 是 $A$ 中相对于
$R$ 的严格标准元；
\item[(f)] (4) $\Rightarrow$ 当 $e \geq e_0$ 时，$a^e$ 是 $A$ 中相对于
$R$ 的初等标准元。
\end{enumerate}
\end{lemma}

\begin{proof}
由定义和《代数》引理 \ref{algebra-lemma-standard-smooth}，(a) 显然。
由定义 \ref{smoothing-definition-strictly-standard}，(b) 显然。

\medskip\noindent
证明 (c)。选取表示 $A = R[x_1, \ldots, x_n]/(f_1, \ldots, f_m)$，
使 (\ref{smoothing-equation-elementary-standard-one}) 和
(\ref{smoothing-equation-elementary-standard-two}) 成立。
选取映到 $a$ 的 $h \in R[x_1, \ldots, x_n]$。于是
$$
A_a = R[x_0, x_1, \ldots, x_n]/(x_0h - 1, f_1, \ldots, f_m).
$$
记 $J = (x_0h - 1, f_1, \ldots, f_m)$。由
(\ref{smoothing-equation-elementary-standard-two})，$A_a$-模 $J/J^2$ 由
$x_0h - 1, f_1, \ldots, f_c$ 生成。因此，与《代数》引理
\ref{algebra-lemma-huber} 的证明相同，可以选取 $g \in 1 + J$，使得
$$
A_a = R[x_0, \ldots, x_n, x_{n + 1}]/
(x_0h - 1, f_1, \ldots, f_c, gx_{n + 1} - 1).
$$
此时，(\ref{smoothing-equation-elementary-standard-one}) 蕴含 $R \to A_a$
标准光滑（用坐标 $x_0, x_1, \ldots, x_c, x_{n + 1}$ 求偏导）。

\medskip\noindent
证明 (d)。选取表示 $A = R[x_1, \ldots, x_n]/(f_1, \ldots, f_m)$，
使 (\ref{smoothing-equation-strictly-standard-one}) 和
(\ref{smoothing-equation-strictly-standard-two}) 成立。记 $I = (f_1, \ldots, f_m)$。
由引理 \ref{smoothing-lemma-elkik}，我们已经知道 $A_a$ 在 $R$ 上光滑。
由引理 \ref{smoothing-lemma-parse-equation-strictly-standard-one}，
$(I/I^2)_a$ 是以 $f_1, \ldots, f_c$ 为基的自由模，并同构地映到
$\bigoplus A_a \text{d}x_i$ 的一个直和项。由于
$\Omega_{A_a/R} = (\Omega_{A/R})_a$ 是同态
$(I/I^2)_a \to \bigoplus A_a \text{d}x_i$ 的余核，故它稳定自由。

\medskip\noindent
证明 (e)。选取表示 $A = R[x_1, \ldots, x_n]/I$，其中 $I$ 有限生成。
由假设，存在短正合列
$$
0 \to (I/I^2)_a \to \bigoplus\nolimits_{i = 1, \ldots, n} A_a\text{d}x_i \to
\Omega_{A_a/R} \to 0
$$
且它分裂正合。因此 $(I/I^2)_a \oplus \Omega_{A_a/R}$ 是自由
$A_a$-模。由于 $\Omega_{A_a/R}$ 稳定自由，$(I/I^2)_a$ 也稳定自由。
于是，对某个 $r$，用表示
$A = R[x_1, \ldots, x_n, x_{n + 1}, \ldots, x_{n + r}]/J$ 代替上述表示，
其中 $J = (I, x_{n + 1}, \ldots, x_{n + r})$，便可使 $(J/J^2)_a$
为（有限）自由模。选取 $f_1, \ldots, f_c \in J$，使其像构成
$(J/J^2)_a$ 的一组基，再把它们扩充为一列生成元
$f_1, \ldots, f_m \in J$。考虑表示
$A = R[x_1, \ldots, x_{n + r}]/(f_1, \ldots, f_m)$。
由构造，当 $e$ 充分大时，$a^e$ 满足
(\ref{smoothing-equation-strictly-standard-two})。此外，由于
$(J/J^2)_a \to \bigoplus\nolimits_{i = 1, \ldots, n + r} A_a\text{d}x_i$
是分裂单射，可以找到一个 $A_a$-线性左逆。用基
$f_1, \ldots, f_c$ 写出此左逆并清去分母，得到线性同态
$\psi_0 : A^{\oplus n + r} \to A^{\oplus c}$，使得
$$
A^{\oplus c} \xrightarrow{(f_1, \ldots, f_c)}
J/J^2 \xrightarrow{f \mapsto \text{d}f}
\bigoplus\nolimits_{i = 1, \ldots, n + r} A \text{d}x_i
\xrightarrow{\psi_0}
A^{\oplus c}
$$
是乘以 $a^{e_0}$ 的同态，其中 $e_0 \geq 1$。由引理
\ref{smoothing-lemma-parse-equation-strictly-standard-one}，
(\ref{smoothing-equation-strictly-standard-one}) 对 $a^{ce_0}$ 成立，
因而对所有满足 $e \geq ce_0$ 的 $a^e$ 成立。

\medskip\noindent
证明 (f)。选取表示
$A_a = R[x_1, \ldots, x_n]/(f_1, \ldots, f_c)$，使得
$\det(\partial f_j/\partial x_i)_{i, j = 1, \ldots, c}$ 在 $A_a$ 中可逆。
由引理 \ref{smoothing-lemma-standard-smooth-include-generators}，可以假设对某个
$m < n$，元素 $x_1, \ldots, x_m$ 的类分别对应于 $a_i/1$，其中
$a_1, \ldots, a_m \in A$ 是 $A$ 在 $R$ 上的一组生成元。
对 $m < i \leq n$ 用 $a^Nx_i$ 代替 $x_i$ 后，可以假设 $x_i$ 的类
为某个 $a_i \in A$ 所给出的 $a_i/1 \in A_a$。考虑环同态
$$
\Psi : R[x_1, \ldots, x_n] \longrightarrow A,\quad
x_i \longmapsto a_i.
$$
这是满环同态。用 $a^Nf_j$ 代替 $f_j$ 后，可以假设
$f_j \in R[x_1, \ldots, x_n]$ 且 $\Psi(f_j) = 0$
（毕竟在 $A_a$ 中 $f_j(a_1/1, \ldots, a_n/1) = 0$）。令
$J = \Ker(\Psi)$。则 $A = R[x_1, \ldots, x_n]/J$ 是一个表示，
并且 $f_1, \ldots, f_c \in J$ 满足：$(J/J^2)_a$ 以
$f_1, \ldots, f_c$ 为自由基，而且
$\det(\partial f_j/\partial x_i)_{i, j = 1, \ldots, c}$ 的像在 $A_a$
中可逆。因此，对所有充分大的 $e$，$a^e$ 满足
(\ref{smoothing-equation-elementary-standard-one}) 和
(\ref{smoothing-equation-elementary-standard-two})，正是所需结论。
\end{proof}






\section{插曲：N\'eron 消奇异化}
\label{smoothing-section-neron}

\noindent
我们暂且中断对 Popescu 定理一般情形的攻克，转而研究一个较容易但已十分
有趣的情形，即 $R \to \Lambda$ 是离散赋值环之间的同态。
对此的讨论见 \cite[Section 4]{Artin-Algebraic-Approximation}。

\begin{situation}
\label{smoothing-situation-neron}
设 $R \subset \Lambda$ 是分歧指数为 $1$ 的离散赋值环扩张
（见《代数详论》定义
\ref{more-algebra-definition-extension-discrete-valuation-rings}）。
给定分解
$$
R \to A \xrightarrow{\varphi} \Lambda
$$
其中 $R \to A$ 平坦且为有限型。令 $\mathfrak q = \Ker(\varphi)$，
$\mathfrak p = \varphi^{-1}(\mathfrak m_\Lambda)$。
\end{situation}

\noindent
在情形 \ref{smoothing-situation-neron} 中，令 $\pi \in R$ 为一致化元。
回顾一下，$A$ 在 $R$ 上平坦意味着 $\pi$ 是 $A$ 上的非零因子
（见《代数详论》引理
\ref{more-algebra-lemma-valuation-ring-torsion-free-flat}）。
由关于 $R \subset \Lambda$ 的假设，$\pi$ 在 $\Lambda$ 中的像也是一致化元。
由于 $\pi \in \mathfrak p$，可以考虑 N\'eron 仿射吹起代数
（见《代数》第 \ref{algebra-section-blow-up} 节）
$$
\varphi' :
A' = A[\textstyle{\frac{\mathfrak p}{\pi}}]
\longrightarrow
\Lambda
$$
它带有到 $\Lambda$ 的诱导同态，把 $a/\pi^n$（其中
$a \in \mathfrak p^n$）映到 $\Lambda$ 中的 $\pi^{-n}\varphi(a)$。
以 $\mathfrak q' \subset \mathfrak p' \subset A'$ 记 $A'$ 中相应的素理想。
注意，$A'$ 的同构类不依赖于一致化元的选取。重复此构造得到序列
$$
A \to A' \to A'' \to \ldots \to \Lambda
$$

\begin{lemma}
\label{smoothing-lemma-neron-functorial}
在情形 \ref{smoothing-situation-neron} 中，N\'eron 吹起在下述意义下具有函子性：
\begin{enumerate}
\item 若 $a \in A$ 且 $a \not \in \mathfrak p$，则 $A_a$ 的 N\'eron 吹起
为 $A'_a$；
\item 若 $B \to A$ 是核为 $I$ 的平坦有限型 $R$-代数之间的满同态，
则 $A'$ 是 $B'/IB'$ 对其 $\pi$-幂挠子模取商所得的环。
\end{enumerate}
\end{lemma}

\begin{proof}
(1) 和 (2) 都是《代数》引理
\ref{algebra-lemma-blowup-base-change} 的特例。事实上，只要有
$A_1 \to A_2 \to \Lambda$ 且 $\mathfrak p_1 A_2 = \mathfrak p_2$，
$A_2'$ 就是 $A_1' \otimes_{A_1} A_2$ 对其 $\pi$-幂挠子模取商所得的环。
\end{proof}

\begin{lemma}
\label{smoothing-lemma-neron-blowup-smooth}
在情形 \ref{smoothing-situation-neron} 中，假设 $R \to A$ 在 $\mathfrak p$ 处光滑，
且 $R/\pi R \subset \Lambda/\pi \Lambda$ 是可分域扩张。
则 $R \to A'$ 在 $\mathfrak p'$ 处光滑，并且存在短正合列
$$
0 \to
\Omega_{A/R} \otimes_A A'_{\mathfrak p'} \to
\Omega_{A'/R, \mathfrak p'} \to
(A'/\pi A')_{\mathfrak p'}^{\oplus c} \to 0
$$
其中 $c = \dim((A/\pi A)_\mathfrak p)$。
\end{lemma}

\begin{proof}
由引理 \ref{smoothing-lemma-neron-functorial}，可以用 $A$ 关于某个不属于
$\mathfrak p$ 的元素的局部化代替 $A$；下文将不再特别说明这一点。
记 $\kappa = R/\pi R$。由于光滑性在基变换下保持
（见《代数》引理 \ref{algebra-lemma-base-change-smooth}），
$A/\pi A$ 在 $\mathfrak p$ 处是光滑 $\kappa$-代数。因此
$(A/\pi A)_\mathfrak p$ 是正则局部环
（见《代数》引理 \ref{algebra-lemma-characterize-smooth-over-field}）。
选取 $g_1, \ldots, g_c \in \mathfrak p$，使其像构成
$(A/\pi A)_\mathfrak p$ 的一组正则参数。必要时局部化 $A$ 后，有
$\mathfrak p = (\pi, g_1, \ldots, g_c)$。注意
$\pi, g_1, \ldots, g_c$ 是 $A_\mathfrak p$ 中的正则序列
（首先 $\pi$ 是非零因子，再对序列的其余部分应用《代数》引理
\ref{algebra-lemma-regular-ring-CM}）。再次局部化 $A$ 后，可以假设
$\pi, g_1, \ldots, g_c$ 是 $A$ 中的正则序列
（见《代数》引理 \ref{algebra-lemma-regular-sequence-in-neighbourhood}）。
由此
$$
A' = A[y_1, \ldots, y_c]/(\pi y_1 - g_1, \ldots, \pi y_c - g_c) =
A[y_1, \ldots, y_c]/I
$$
由《代数详论》引理 \ref{more-algebra-lemma-blowup-regular-sequence} 得出。
下文将直接使用以朴素余切复形表述的光滑性定义
（见《代数》定义 \ref{algebra-definition-smooth}）以及《代数》引理
\ref{algebra-lemma-smooth-at-point} 的判别准则，不再另行说明。
《代数》引理 \ref{algebra-lemma-exact-sequence-NL} 对
$R \to A[y_1, \ldots, y_c] \to A'$ 给出的正合列为
$$
0 \to H_1(\NL_{A'/R}) \to I/I^2 \to
\Omega_{A/R} \otimes_A A' \oplus
\bigoplus\nolimits_{i = 1, \ldots, c} A' \text{d}y_i \to
\Omega_{A'/R} \to 0
$$
其中 $I/I^2$ 中 $\pi y_i - g_i$ 的类映到下一项中的
$- \text{d}g_i + \pi \text{d}y_i$。这里用《代数》引理
\ref{algebra-lemma-NL-surjection} 计算
$\NL_{A'/A[y_1, \ldots, y_c]}$，并使用了
$R \to A[y_1, \ldots, y_c]$ 光滑这一事实。因此
$H_1(\NL_{A[y_1, \ldots, y_c]/R}) = 0$，而且
$\Omega_{A[y_1, \ldots, y_c]/R}$ 是有限投射（因而平坦）
$A[y_1, \ldots, y_c]$-模；事实上，它是
$\Omega_{A/R} \otimes_A A[y_1, \ldots, y_c]$ 与以 $\text{d}y_i$
为基的自由模之直和。为完成证明，只需说明
$\text{d}g_1, \ldots, \text{d}g_c$ 是有限自由模
$\Omega_{A/R, \mathfrak p}$ 的一组基的一部分。事实上，这将说明
$(I/I^2)_\mathfrak p$ 是以 $\pi y_i - g_i$ 为基的自由模，
正合列中间同态在 $\mathfrak p$ 处的局部化是单射，因而
$H_1(\NL_{A'/R})_\mathfrak p = 0$，并且余核
$\Omega_{A'/R, \mathfrak p}$ 有限自由。为此，只需证明
$\text{d}g_i$ 的像在
$\Omega_{A/R, \mathfrak p}/\pi = \Omega_{(A/\pi A)/\kappa, \mathfrak p}$
中在 $\kappa(\mathfrak p)$ 上线性无关（该等式由《代数》引理
\ref{algebra-lemma-differentials-base-change} 给出）。由于
$\kappa \subset \kappa(\mathfrak p) \subset \Lambda/\pi \Lambda$，
$\kappa(\mathfrak p)$ 在 $\kappa$ 上可分
（见《代数》定义 \ref{algebra-definition-separable-field-extension}）。
所需的线性无关性由《代数》引理
\ref{algebra-lemma-computation-differential} 得出。
\end{proof}

\begin{lemma}
\label{smoothing-lemma-neron-when-smooth}
在情形 \ref{smoothing-situation-neron} 中，假设 $R \to A$ 在 $\mathfrak q$ 处光滑，
并有核为 $I$ 的满 $R$-代数同态 $B \to A$。假设 $R \to B$ 在
$\mathfrak p_B = (B \to A)^{-1}\mathfrak p$ 处光滑。若同态
$$
I/I^2 \otimes_A \Lambda \to \Omega_{B/R} \otimes_B \Lambda
$$
的余核为自由 $\Lambda$-模，则 $R \to A$ 在 $\mathfrak p$ 处光滑。
\end{lemma}

\begin{proof}
同态 $I/I^2 \to \Omega_{B/R} \otimes_B A$ 的余核为 $\Omega_{A/R}$，
见《代数》引理 \ref{algebra-lemma-differential-seq}。令
$d = \dim_\mathfrak q(A/R)$ 为 $R \to A$ 在 $\mathfrak q$ 处的相对维数，
即 $\Spec(A[1/\pi])$ 在 $\mathfrak q$ 处的维数；见《代数》定义
\ref{algebra-definition-relative-dimension}。于是
$\Omega_{A/R, \mathfrak q}$ 是秩为 $d$ 的自由 $A_\mathfrak q$-模
（见《代数》引理 \ref{algebra-lemma-characterize-smooth-over-field}）。
因此，若引理的假设成立，则 $\Omega_{A/R} \otimes_A \Lambda$
是秩为 $d$ 的自由模。从而 $\Omega_{A/R} \otimes_A \kappa(\mathfrak p)$
的维数为 $d$（与 $\Lambda/\pi \Lambda$ 张量后即可看出）。
由于 $R \to A$ 平坦，且 $\mathfrak p$ 是 $\mathfrak q$ 的特化，
《代数》引理 \ref{algebra-lemma-dimension-fibres-bounded-open-upstairs}
给出 $\dim_\mathfrak p(A/R) \geq d$。最后，由《代数》引理
\ref{algebra-lemma-flat-fibre-smooth} 和
\ref{algebra-lemma-characterize-smooth-over-field}，
$R \to A$ 在 $\mathfrak p$ 处光滑。
\end{proof}

\begin{lemma}
\label{smoothing-lemma-neron-desingularization}
在情形 \ref{smoothing-situation-neron} 中，假设 $R \to A$ 在 $\mathfrak q$ 处光滑，
且 $R/\pi R \subset \Lambda/\pi \Lambda$ 是可分域扩张。
则经过有限次仿射 N\'eron 吹起后，代数 $A$ 在 $\mathfrak p$ 处变为
光滑 $R$-代数。
\end{lemma}

\begin{proof}
选取 $R$-代数 $B$ 和满同态 $B \to A$。令
$\mathfrak p_B = (B \to A)^{-1}(\mathfrak p)$，并以 $r$ 记
$R \to B$ 在 $\mathfrak p_B$ 处的相对维数。选取 $B$，使得
$R \to B$ 在 $\mathfrak p_B$ 处光滑。例如，可取 $B$ 为 $R$ 上有
$r$ 个变量的多项式代数。考虑引理 \ref{smoothing-lemma-neron-when-smooth} 中的复形
$$
I/I^2 \otimes_A \Lambda \longrightarrow \Omega_{B/R} \otimes_B \Lambda
$$
由有限 $\Lambda$-模的结构
（见《代数详论》引理 \ref{more-algebra-lemma-modules-PID}），
其余核形如
$$
\Lambda^{\oplus d} \oplus
\bigoplus\nolimits_{i = 1, \ldots, n} \Lambda/\pi^{e_i} \Lambda
$$
其中 $d \geq 0$、$n \geq 0$ 且 $e_i \geq 1$。注意，$d$ 是 $A/R$
在 $\mathfrak q$ 处的相对维数
（见《代数》引理 \ref{algebra-lemma-characterize-smooth-over-field}）。
若缺陷 $e = \sum_{i = 1, \ldots, n} e_i$ 为零，
则引理 \ref{smoothing-lemma-neron-when-smooth} 已给出结论。

\medskip\noindent
下面考察施行 N\'eron 吹起后发生的情形。回顾一下，$A'$ 是
$B'/IB'$ 对其 $\pi$-幂挠子模取商所得的环
（见引理 \ref{smoothing-lemma-neron-functorial}），而 $R \to B'$ 在
$\mathfrak p_{B'}$ 处光滑（见引理 \ref{smoothing-lemma-neron-blowup-smooth}）。
因此吹起后得到完全相同的设定。图示如下：
$$
\xymatrix{
0 \ar[r] & I' \ar[r] & B' \ar[r] & A' \ar[r] & 0 \\
0 \ar[r] & I \ar[u] \ar[r] & B \ar[u] \ar[r] & A \ar[r] \ar[u] & 0
}
$$
由于 $I \subset \mathfrak p_B$，同态 $I \to I'$ 经由 $\pi I'$ 分解。
考察所诱导的复形同态，得到
$$
\xymatrix{
I'/(I')^2 \otimes_{A'} \Lambda \ar[r] &
\Omega_{B'/R} \otimes_{B'} \Lambda \ar@{=}[r] & M' \\
I/I^2 \otimes_A \Lambda \ar[r] \ar[u] &
\Omega_{B/R} \otimes_B \Lambda \ar[u] \ar@{=}[r] & M
}
$$
此时 $M \subset M'$ 是有限自由 $\Lambda$-模，且商 $M'/M$ 被 $\pi$ 零化，
见引理 \ref{smoothing-lemma-neron-blowup-smooth}。令 $N \subset M$ 与
$N' \subset M'$ 分别为水平同态的像，并记 $Q = M/N$、$Q' = M'/N'$。
得到交换图
$$
\xymatrix{
0 \ar[r] &
N' \ar[r] &
M' \ar[r] &
Q' \ar[r] &
0 \\
0 \ar[r] &
N \ar[r] \ar[u] &
M \ar[r] \ar[u] &
Q \ar[r] \ar[u] &
0
}
$$
其中 $N \subset N'$ 是秩为 $r - d$ 的自由 $\Lambda$-模。
由于 $I$ 映入 $\pi I'$，有 $N \subset \pi N'$。

\medskip\noindent
令 $K = \Lambda_\pi$ 为 $\Lambda$ 的分式域。存在交换图
$$
\xymatrix{
0 \ar[r] &
N' \ar[r] &
N'_K \cap M' \ar[r] &
Q'_{tor} \ar[r] &
0 \\
0 \ar[r] &
N \ar[r] \ar[u] &
N_K \cap M \ar[r] \ar[u] &
Q_{tor} \ar[r] \ar[u] &
0
}
$$
其两行都是短正合列。这表明缺陷的变化量为
$$
e - e' =
\text{length}(Q_{tor}) - \text{length}(Q'_{tor})
=
\text{length}(N'/N) - \text{length}(N'_K \cap M' / N_K \cap M)
$$
由于 $M'/M$ 被 $\pi$ 零化，$N'_K \cap M' / N_K \cap M$ 也被
$\pi$ 零化，故其长度至多为 $\dim_K(N_K)$。又因
$N \subset \pi N'$，有 $\text{length}(N'/N) \ge \dim_K(N_K)$，
而等号成立当且仅当 $N = \pi N'$。

\medskip\noindent
为完成证明，需要说明：当 $A$ 在 $\mathfrak p$ 处不光滑时，
$N$ 严格小于 $\pi N'$；这正是 N\'eron 论证中的关键计算。
为此，考虑正合列
$$
I/I^2 \otimes_B \kappa(\mathfrak p_B)
\to \Omega_{B/R} \otimes_B \kappa(\mathfrak p_B)
\to \Omega_{A/R} \otimes_A \kappa(\mathfrak p) \to 0
$$
（由《代数》引理 \ref{algebra-lemma-differential-seq} 得出）。
由于 $R \to A$ 在 $\mathfrak p$ 处不光滑，
$\Omega_{A/R} \otimes_A \kappa(\mathfrak p)$ 的维数 $s$ 大于 $d$。
另一方面，第一个同态经由单射
$$
\mathfrak p B_\mathfrak p/\mathfrak p^2 B_\mathfrak p
\to \Omega_{B/R} \otimes_B \kappa(\mathfrak p_B)
$$
分解；该单射来自《代数》引理
\ref{algebra-lemma-computation-differential}。注意，由关于
$R/\pi R \subset \Lambda/\pi \Lambda$ 的假设，
$\kappa(\mathfrak p)$ 在 $k$ 上可分。因此可以找到生成元
$g_1, \ldots, g_t \in I$，使得当 $j > r - s$ 时
$g_j \in \mathfrak p^2$。于是当 $j > r - s$ 时，$g_j$ 在 $A'$ 中的像
属于 $\pi^2 I'$。由于 $r - s < r - d$，$N$ 的至少一个极小生成元
在 $N'$ 中变得可被 $\pi^2$ 整除。因此缺陷 $e$ 至少减少 $1$，证明完成。
\end{proof}

\noindent
若 $R \to \Lambda$ 是离散赋值环扩张，则 $R \to \Lambda$ 正则当且仅当：
(a) 分歧指数为 $1$；(b) 分式域扩张可分；并且
(c) $R/\mathfrak m_R \subset \Lambda/\mathfrak m_\Lambda$ 可分。
因此，以下结果是定理 \ref{smoothing-theorem-popescu} 中一般 N\'eron 消奇异化的特例。

\begin{lemma}
\label{smoothing-lemma-neron-colimit}
\begin{slogan}
DVR 的非分歧扩张是 ind-smooth 的，亦即 N\'eron 消奇异化
\end{slogan}
设 $R \subset \Lambda$ 是分歧指数为 $1$ 的离散赋值环扩张，
并且它所诱导的剩余域扩张和分式域扩张均可分。
则 $\Lambda$ 是光滑 $R$-代数的滤过余极限。
\end{lemma}

\begin{proof}
由《代数》引理 \ref{algebra-lemma-when-colimit}，只需证明情形
\ref{smoothing-situation-neron} 中的任意 $R \to A \to \Lambda$ 都可分解为
$A \to B \to \Lambda$，其中 $B$ 是光滑 $R$-代数。用 $A$ 在
$\Lambda$ 中的像代替 $A$ 后，可以假设 $A$ 是整环，其分式域 $K$
是 $\Lambda$ 的分式域的子域。特别地，由假设，$A$ 在 $R$ 的分式域上可分。
于是由《代数》引理 \ref{algebra-lemma-smooth-at-generic-point}，
$R \to A$ 在 $\mathfrak q = (0)$ 处光滑。经有限次 N\'eron 吹起后，
可以假设 $R \to A$ 在 $\mathfrak p$ 处光滑，见引理
\ref{smoothing-lemma-neron-desingularization}。随后，对某个
$a \in A$、$a \not \in \mathfrak p$，用其局部化代替 $A$，即可使
$A$ 在 $R$ 上光滑。引理得证。
\end{proof}












\section{提升问题}
\label{smoothing-section-lifting}

\noindent
本节的目标是证明（命题 \ref{smoothing-proposition-lift}）：由光滑代数的滤过余极限
构成的代数类在平坦无穷小形变下封闭。证明是初等的，只用到第
\ref{smoothing-section-presentations} 节关于光滑代数表示的结果。

\begin{lemma}
\label{smoothing-lemma-lift-once}
设 $R \to \Lambda$ 为环同态，$I \subset R$ 为理想。假设
\begin{enumerate}
\item $I^2 = 0$；
\item $\Lambda/I\Lambda$ 是光滑 $R/I$-代数的滤过余极限。
\end{enumerate}
设 $\varphi : A \to \Lambda$ 为 $R$-代数同态，其中 $A$ 是有限表示
$R$-代数。则存在分解
$$
A \to B/J \to \Lambda
$$
其中 $B$ 是光滑 $R$-代数，而 $J \subset IB$ 是有限生成理想。
\end{lemma}

\begin{proof}
选取分解
$$
A/IA \to \bar B \to \Lambda/I\Lambda
$$
其中 $\bar B$ 在 $R/I$ 上标准光滑；由假设和引理
\ref{smoothing-lemma-colimit-standard-smooth}，这样的分解存在。写成
$$
\bar B = A/IA[t_1, \ldots, t_r]/(\bar g_1, \ldots, \bar g_s)
$$
并设 $\bar B \to \Lambda/I\Lambda$ 把 $t_i$ 映到 $\lambda_i$ 模
$I\Lambda$ 的类。选取 $\bar g_1, \ldots, \bar g_s$ 的提升
$g_1, \ldots, g_s \in A[t_1, \ldots, t_r]$。写成
$\varphi(g_i)(\lambda_1, \ldots, \lambda_r) =
\sum \epsilon_{ij} \mu_{ij}$
其中某些 $\epsilon_{ij} \in I$，$\mu_{ij} \in \Lambda$。定义
$$
A' = A[t_1, \ldots, t_r, \delta_{i, j}]/
(g_i - \sum \epsilon_{ij} \delta_{ij})
$$
并考虑同态
$$
A' \longrightarrow \Lambda,\quad
a \longmapsto \varphi(a),\quad
t_i \longmapsto \lambda_i,\quad
\delta_{ij} \longmapsto \mu_{ij}
$$
有
$$
A'/IA' = A/IA[t_1, \ldots, t_r]/(\bar g_1, \ldots, \bar g_s)[\delta_{ij}]
\cong \bar B[\delta_{ij}]
$$
由于 $\bar B$ 标准光滑，这是一个标准光滑 $R/I$-代数。选取表示
$A'/IA' = R/I[x_1, \ldots, x_n]/(\bar f_1, \ldots, \bar f_c)$，使得
$\det(\partial \bar f_j/\partial x_i)_{i, j = 1, \ldots, c}$ 在
$A'/IA'$ 中可逆。选取 $\bar f_1, \ldots, \bar f_c$ 的提升
$f_1, \ldots, f_c \in R[x_1, \ldots, x_n]$。于是
$$
B = R[x_1, \ldots, x_n, x_{n + 1}]/
(f_1, \ldots, f_c,
x_{n + 1}\det(\partial f_j/\partial x_i)_{i, j = 1, \ldots, c} - 1)
$$
在 $R$ 上光滑。由于光滑环同态形式光滑
（见《代数》命题 \ref{algebra-proposition-smooth-formally-smooth}），
存在模 $I$ 后为同构的 $R$-代数同态 $B \to A'$。再由 Nakayama 引理
（《代数》引理 \ref{algebra-lemma-NAK}），$B \to A'$ 是满射。
因此 $A' = B/J$，其中 $J \subset IB$ 有限生成
（见《代数》引理 \ref{algebra-lemma-finite-presentation-independent}）。
\end{proof}

\begin{lemma}
\label{smoothing-lemma-lift-twice}
设 $R \to \Lambda$ 为环同态，$I \subset R$ 为理想。假设
\begin{enumerate}
\item $I^2 = 0$；
\item $\Lambda/I\Lambda$ 是光滑 $R/I$-代数的滤过余极限；
\item $R \to \Lambda$ 平坦。
\end{enumerate}
设 $\varphi : B \to \Lambda$ 为 $R$-代数同态，其中 $B$ 在 $R$ 上光滑。
设 $J \subset IB$ 为有限生成理想，且 $\varphi(J) = 0$。
则存在 $R$-代数同态
$$
B \xrightarrow{\alpha} B' \xrightarrow{\beta} \Lambda
$$
使得 $B'$ 在 $R$ 上光滑、$\alpha(J) = 0$，并且
$\beta \circ \alpha = \varphi$。
\end{lemma}

\begin{proof}
只要能对 $J = (h)$ 的情形证明引理，便可对 $J$ 的生成元个数作归纳。
具体地，假设 $J$ 可由 $n$ 个元素 $h_1, \ldots, h_n$ 生成，且引理对
$J$ 由 $n - 1$ 个元素生成的一切情形成立。先应用 $n = 1$ 的情形，
得到 $B \to B' \to \Lambda$，其中第一个同态把 $h_n$ 映为零。
再令 $J'$ 为 $B'$ 中由 $h_1, \ldots, h_{n - 1}$ 的像生成的理想，
应用 $n - 1$ 的情形得到 $B' \to B'' \to \Lambda$。
容易验证 $B \to B'' \to \Lambda$ 满足要求。

\medskip\noindent
假设 $J = (h)$，并把 $h$ 写成 $h = \sum \epsilon_i b_i$，其中
$\epsilon_i \in I$、$b_i \in B$。注意
$0 = \varphi(h) = \sum \epsilon_i \varphi(b_i)$。由于 $\Lambda$ 在
$R$ 上平坦，平坦性的方程判别准则
（见《代数》引理 \ref{algebra-lemma-flat-eq}）表明，存在
$\lambda_j \in \Lambda$（$j = 1, \ldots, m$）和 $a_{ij} \in R$，使得
$\varphi(b_i) = \sum_j a_{ij} \lambda_j$ 且
$\sum_i \epsilon_i a_{ij} = 0$。令
$$
C = B[x_1, \ldots, x_m]/(b_i - \sum a_{ij} x_j)
$$
并以 $\varphi$ 和 $x_j \mapsto \lambda_j$ 定义 $C \to \Lambda$。
按引理 \ref{smoothing-lemma-lift-once} 选取分解
$$
C \to B'/J' \to \Lambda
$$
由于 $B$ 在 $R$ 上光滑，可以把同态 $B \to C \to B'/J'$ 提升为
$\alpha : B \to B'$。于是 $\varphi = \beta \circ \alpha$。
为完成证明，只需验证 $\alpha(h) = 0$。事实上，$\alpha$ 是
$B \to C \to B'/J'$ 的提升这一事实蕴含
$$
\alpha(b_i) = \sum a_{ij} \xi_j + \theta_i
$$
其中某些 $\xi_j \in B'$，$\theta_i \in J' \subset IB'$。因此
$$
\alpha(h) = \alpha(\sum \epsilon_i b_i) =
\sum \epsilon_i a_{ij} \xi_j + \sum \epsilon_i \theta_i = 0
$$
这是因为上述关系成立且 $I^2 = 0$。
\end{proof}

\begin{proposition}
\label{smoothing-proposition-lift}
\begin{slogan}
代数的 ind-光滑性在无穷小形变下保持
\end{slogan}
设 $R \to \Lambda$ 为环同态，$I \subset R$ 为理想。假设
\begin{enumerate}
\item $I$ 是幂零理想；
\item $\Lambda/I\Lambda$ 是光滑 $R/I$-代数的滤过余极限；
\item $R \to \Lambda$ 平坦。
\end{enumerate}
则 $\Lambda$ 是光滑 $R$-代数的滤过余极限。
\end{proposition}

\begin{proof}
对某个 $n$ 有 $I^n = 0$，故对 $n$ 作归纳可知，只需考虑 $I^2 = 0$
的情形。设 $\varphi : A \to \Lambda$ 为 $R$-代数同态，其中 $A$ 是
有限表示 $R$-代数。由《代数》引理 \ref{algebra-lemma-when-colimit}，
需要找到分解 $A \to B \to \Lambda$，其中 $B$ 在 $R$ 上光滑。
由引理 \ref{smoothing-lemma-lift-once}，可以假设 $A = B/J$，其中 $B$ 在 $R$
上光滑，而 $J \subset IB$ 是有限生成理想。由引理
\ref{smoothing-lemma-lift-twice}，可以找到交换图
$$
\xymatrix{
B \ar[rr]_\alpha \ar[rd]_\varphi & & B' \ar[ld]^\beta \\
& \Lambda
}
$$
其中各同态均为 $R$-代数同态，$B'$ 在 $R$ 上光滑，且
$\alpha(J) = 0$。所以 $\alpha$ 分解为 $B \to A \to B'$，证明完成。
\end{proof}






\section{提升引理}
\label{smoothing-section-lifting-lemma}

\noindent
下面是一个极其巧妙的引理。

\begin{lemma}
\label{smoothing-lemma-lifting}
设 $R$ 为诺特环，$\Lambda$ 为 $R$-代数。设 $\pi \in R$，并假设
$\text{Ann}_R(\pi) = \text{Ann}_R(\pi^2)$ 且
$\text{Ann}_\Lambda(\pi) = \text{Ann}_\Lambda(\pi^2)$。
假设有 $R$-代数同态
$R/\pi^2R \to \bar C \to \Lambda/\pi^2\Lambda$，其中 $\bar C$ 有限表示。
则存在 $R$-代数同态 $D \to \Lambda$ 及交换图
$$
\xymatrix{
R/\pi^2R \ar[r] \ar[d] &
\bar C \ar[r] \ar[d] &
\Lambda/\pi^2\Lambda \ar[d] \\
R/\pi R \ar[r] &
D/\pi D \ar[r] &
\Lambda/\pi \Lambda
}
$$
满足以下性质：
\begin{enumerate}
\item[(a)] $D$ 有限表示；
\item[(b)] 对任意满足 $\pi \not \in \mathfrak q$ 的素理想 $\mathfrak q$，
$R \to D$ 在 $\mathfrak q$ 处光滑；
\item[(c)] 对任意满足 $\pi \in \mathfrak q$、且位于 $\bar C$ 中一个
$R/\pi^2 R \to \bar C$ 光滑的素理想之上的素理想 $\mathfrak q$，
$R \to D$ 在 $\mathfrak q$ 处光滑；
\item[(d)] 对任意位于 $\bar C$ 中一个 $R/\pi^2R \to \bar C$ 光滑的
素理想之上的素理想，$\bar C/\pi \bar C \to D/\pi D$ 在该处光滑。
\end{enumerate}
\end{lemma}

\begin{proof}
选取表示
$$
\bar C = R[x_1, \ldots, x_n]/(f_1, \ldots, f_m)
$$
再记 $I = (f_1, \ldots, f_m)$，并以 $\bar I$ 记 $I$ 在
$R/\pi^2R[x_1, \ldots, x_n]$ 中的像。由于 $R$ 是诺特环，$\bar C$
也是诺特环。因此 $R/\pi^2 R \to \bar C$ 的光滑轨迹拟紧，见
《拓扑》引理 \ref{topology-lemma-Noetherian}。应用引理
\ref{smoothing-lemma-find-strictly-standard}，可以选取有限个元素
$a_1, \ldots, a_r \in R[x_1, \ldots, x_n]$，使得
\begin{enumerate}
\item 开子空间 $\Spec(\bar C_{a_k}) \subset \Spec(\bar C)$ 的并覆盖
$R/\pi^2 R \to \bar C$ 的光滑轨迹；
\item 对每个 $k = 1, \ldots, r$，存在有限子集
$E_k \subset \{1, \ldots, m\}$，使得 $(\bar I/\bar I^2)_{a_k}$
以 $f_j$（$j \in E_k$）的类为自由基。
\end{enumerate}
令 $I_k = (f_j, j \in E_k) \subset I$，并以 $\bar I_k$ 记 $I_k$ 在
$R/\pi^2R[x_1, \ldots, x_n]$ 中的像。由 (2) 和 Nakayama 引理，
对某个 $b'_k \in \bar I_{a_k}$，$(\bar I/\bar I_k)_{a_k}$ 被
$1 + b'_k$ 零化。设 $b'_k$ 是某个 $b_k \in I$ 和某个整数 $N$ 所给出的
$b_k/(a_k)^N$ 的像。用 $a_kb_k$ 代替 $a_k$ 后，得到
\begin{enumerate}
\item[(3)] $(\bar I_k)_{a_k} = (\bar I)_{a_k}$.
\end{enumerate}
因此，必要时用 $a_k$ 的一个高次幂代替 $a_k$ 后，可以写成
\begin{enumerate}
\item[(4)]
$a_k f_\ell = \sum\nolimits_{j \in E_k} h_{k, \ell}^jf_j + \pi^2 g_{k, \ell}$
\end{enumerate}
其中 $\ell \in \{1, \ldots, m\}$ 任意，而某些
$h_{i, \ell}^j, g_{i, \ell} \in R[x_1, \ldots, x_n]$。
若 $\ell \in E_k$，则取 $h_{k, \ell}^j = a_k\delta_{\ell, j}$
（Kronecker delta）和 $g_{k, \ell} = 0$。令
$$
D = R[x_1, \ldots, x_n, z_1, \ldots, z_m]/
(f_j - \pi z_j, p_{k, \ell}).
$$
这里 $j \in \{1, \ldots, m\}$、$k \in \{1, \ldots, r\}$、
$\ell \in \{1, \ldots, m\}$，并且
$$
p_{k, \ell} = a_k z_\ell - \sum\nolimits_{j \in E_k} h_{k, \ell}^j z_j
- \pi g_{k, \ell}.
$$
注意，由上述选取，当 $\ell \in E_k$ 时 $p_{k, \ell} = 0$。

\medskip\noindent
同态 $R \to D$ 取为给定同态。设 $\bar C \to \Lambda/\pi^2\Lambda$
把 $x_i$ 映到 $\lambda_i$ 模 $\pi^2$ 的类。对元素
$f \in R[x_1, \ldots, x_n]$，以 $f(\lambda) \in \Lambda$ 记把
$\lambda_i$ 代入 $x_i$ 后的结果。于是对某些 $\mu_j \in \Lambda$，有
$f_j(\lambda) = \pi^2 \mu_j$。按 $x_i \mapsto \lambda_i$ 和
$z_j \mapsto \pi\mu_j$ 定义 $D \to \Lambda$。这是良定义的，因为
\begin{align*}
p_{k, \ell} & \mapsto
a_k(\lambda) \pi \mu_\ell -
\sum\nolimits_{j \in E_k} h_{k, \ell}^j(\lambda) \pi \mu_j
- \pi g_{k, \ell}(\lambda) \\
& =
\pi\left(a_k(\lambda) \mu_\ell -
\sum\nolimits_{j \in E_k} h_{k, \ell}^j(\lambda) \mu_j
- g_{k, \ell}(\lambda)\right)
\end{align*}
在上述 (4) 中代入 $x_i = \lambda_i$，可见括号内的表达式被 $\pi^2$
零化；由于假设 $\text{Ann}_\Lambda(\pi) = \text{Ann}_\Lambda(\pi^2)$，
它也被 $\pi$ 零化。同态 $\bar C \to D/\pi D$ 由
$x_i \mapsto x_i$ 决定（显然良定义）。因此，只须证明 (b)、(c) 和 (d)。

\medskip\noindent
利用 (4) 得到以下关键等式：
\begin{align*}
\pi p_{k, \ell} & =
\pi a_k z_\ell - \sum\nolimits_{j \in E_k} \pi h_{k, \ell}^jz_j
- \pi^2 g_{k, \ell} \\
& =
- a_k (f_\ell - \pi z_\ell) + a_k f_\ell +
\sum\nolimits_{j \in E_k} h_{k, \ell}^j (f_j - \pi z_j) -
\sum\nolimits_{j \in E_k} h_{k, \ell}^j f_j - \pi^2 g_{k, \ell} \\
& =
-a_k(f_\ell - \pi z_\ell) +
\sum\nolimits_{j \in E_k} h_{k, \ell}^j(f_j - \pi z_j)
\end{align*}
最后所得表达式属于由 $f_j - \pi z_j$ 生成的理想。特别地，
$D[1/\pi]$ 同构于
$R[1/\pi][x_1, \ldots, x_n, z_1, \ldots, z_m]/(f_j - \pi z_j)$，
后者又同构于 $R[1/\pi][x_1, \ldots, x_n]$，因而在 $R$ 上光滑。
这证明了 (b)。

\medskip\noindent
固定 $k \in \{1, \ldots, r\}$，考虑环
$$
D_k = R[x_1, \ldots, x_n, z_1, \ldots, z_m]/
(f_j - \pi z_j, j \in E_k, p_{k, \ell})
$$
由于当 $\ell \in E_k$ 时 $p_{k, \ell}$ 为零，方程的数目为
$m = |E_k| + (m - |E_k|)$。还要注意
\begin{align*}
(D_k/\pi D_k)_{a_k}
& =
R/\pi R[x_1, \ldots, x_n, 1/a_k, z_1, \ldots, z_m]/
(f_j, j \in E_k, p_{k, \ell}) \\
& =
(\bar C/\pi \bar C)_{a_k}[z_1, \ldots, z_m]/
(a_kz_\ell - \sum\nolimits_{j \in E_k} h_{k, \ell}^j z_j) \\
& \cong
(\bar C/\pi \bar C)_{a_k}[z_j, j \in E_k]
\end{align*}
特别地，$(D_k/\pi D_k)_{a_k}$ 在 $(\bar C/\pi \bar C)_{a_k}$ 上光滑。
由我们对 $a_k$ 的选择，$(\bar C/\pi \bar C)_{a_k}$ 在
$R/\pi R$ 上光滑且相对维数为 $n - |E_k|$，见 (2)。因此，若素理想
$\mathfrak q_k \subset D_k$ 含有 $\pi$ 且位于
$\Spec(\bar C_{a_k})$ 之上，则 $R \to D_k$ 的纤维环
在 $\mathfrak q_k$ 处光滑且维数为 $n$。于是，由上述方程数目的计算，
$R \to D_k$ 在 $\mathfrak q_k$ 处为 syntomic，见《代数》
引理 \ref{algebra-lemma-localize-relative-complete-intersection}。
因而 $R \to D_k$ 在 $\mathfrak q_k$ 处光滑，见《代数》
引理 \ref{algebra-lemma-flat-fibre-smooth}。

\medskip\noindent
为完成证明，设 $\mathfrak q \subset D$ 是一个含有 $\pi$ 的素理想，
它位于某个使 $R/\pi^2 R \to \bar C$ 光滑的素理想之上。
由 (1)，对某个 $k$ 有 $a_k \not \in \mathfrak q$。
我们将证明满射 $D_k \to D$ 在 $\mathfrak q$ 处诱导局部环的同构。
由上一段，我们已经知道环同态
$\bar C/\pi \bar C \to D_k/\pi D_k$ 和 $R \to D_k$
在相应的素理想 $\mathfrak q_k$ 处光滑。这将证明 (c) 和 (d)，
从而完成整个证明。

\medskip\noindent
首先注意，对任意 $\ell$，上面证明的等式
$\pi p_{k, \ell} = -a_k(f_\ell - \pi z_\ell) +
\sum_{j \in E_k} h_{k, \ell}^j (f_j - \pi z_j)$ 表明
$f_\ell - \pi z_\ell$ 在 $(D_k)_{a_k}$ 中的像为零，因而它在
$(D_k)_{\mathfrak q_k}$ 中的像尤其为零。
关系式 (4) 蕴含在 $I/I^2$ 中有 $a_k f_\ell =
\sum_{j \in E_k} h_{k, \ell}^j f_j$。
由于 $(\bar I_k/\bar I_k^2)_{a_k}$ 是以
$f_j$（$j \in E_k$）为基的自由模，可知
$$
a_{k'} h_{k, \ell}^j -
\sum\nolimits_{j' \in E_{k'}} h_{k', \ell}^{j'} h_{k, j'}^j
$$
对每个 $k,k',\ell$ 和 $j \in E_k$，这个元素在 $\bar C_{a_k}$ 中为零。
因此，可以取充分大的整数 $N$，使得
$$
a_k^N\left(
a_{k'} h_{k, \ell}^j -
\sum\nolimits_{j' \in E_{k'}} h_{k', \ell}^{j'} h_{k, j'}^j
\right)
$$
属于 $I_k + \pi^2R[x_1, \ldots, x_n]$。模 $\pi$ 计算可得
\begin{align*}
&
a_kp_{k', \ell} - a_{k'}p_{k, \ell} + \sum h_{k', \ell}^{j'} p_{k, j'}
\\
&
=
- a_k \sum h_{k', \ell}^{j'} z_{j'}
+ a_{k'} \sum h_{k, \ell}^j z_j
+ \sum h_{k', \ell}^{j'} a_k z_{j'}
- \sum \sum h_{k', \ell}^{j'} h_{k, j'}^j z_j \\
&
=
\sum \left(
a_{k'} h_{k, \ell}^j
- \sum h_{k', \ell}^{j'} h_{k, j'}^j
\right) z_j
\end{align*}
其中采用爱因斯坦求和约定。结合前述结果可知，在
$R[x_1, \ldots, x_n, z_1, \ldots, z_m]$ 中，
$a_k^{N + 1} p_{k', \ell}$ 属于由 $I_k$ 和 $\pi$ 生成的理想。
因此 $p_{k', \ell}$ 的像属于 $\pi (D_k)_{a_k}$。另一方面，等式
$$
\pi p_{k', \ell} =
-a_{k'} (f_\ell - \pi z_\ell) +
\sum\nolimits_{j' \in E_{k'}} h_{k', \ell}^{j'}(f_{j'} - \pi z_{j'})
$$
表明 $\pi p_{k', \ell}$ 在 $(D_k)_{a_k}$ 中为零。
我们已经假设 $\text{Ann}_R(\pi) = \text{Ann}_R(\pi^2)$，
而 $(D_k)_{\mathfrak q_k}$ 在 $R$ 上光滑，因而在 $R$ 上平坦，故有
$\text{Ann}_{(D_k)_{\mathfrak q_k}}(\pi) =
\text{Ann}_{(D_k)_{\mathfrak q_k}}(\pi^2)$.
于是 $p_{k', \ell}$ 的像也为零，从而
$D_{\mathfrak q} = (D_k)_{\mathfrak q_k}$，证毕。
\end{proof}






\section{消奇异化引理}
\label{smoothing-section-desingularization-lemma}

\noindent
下面是另一个极为巧妙的引理。

\begin{lemma}
\label{smoothing-lemma-desingularize}
设 $R$ 为诺特环，$\Lambda$ 为 $R$-代数。取 $\pi \in R$，
并假设 $\text{Ann}_\Lambda(\pi) = \text{Ann}_\Lambda(\pi^2)$。
设 $A \to \Lambda$ 为 $R$-代数同态，其中 $A$ 为有限表示。假设
\begin{enumerate}
\item $\pi$ 的像在 $A$ 中相对于 $R$ 是严格标准元；
\item 存在与到 $\Lambda/\pi^4 \Lambda$ 的同态相容的截面
$\rho : A/\pi^4 A \to R/\pi^4 R$。
\end{enumerate}
则可以找到 $R$-代数同态 $A \to B \to \Lambda$，其中 $B$ 为有限表示，
并且 $\mathfrak a B \subset H_{B/R}$，这里
$\mathfrak a = \text{Ann}_R(\text{Ann}_R(\pi^2)/\text{Ann}_R(\pi))$。
\end{lemma}

\begin{proof}
选取一个表示
$$
A = R[x_1, \ldots, x_n]/(f_1, \ldots, f_m)
$$
以及 $0 \leq c \leq \min(n, m)$，使得
(\ref{smoothing-equation-strictly-standard-one}) 对 $\pi$ 成立，并且
\begin{equation}
\label{smoothing-equation-star}
\pi f_{c + j} \in (f_1, \ldots, f_c) + (f_1, \ldots, f_m)^2
\end{equation}
对 $j = 1, \ldots, m - c$ 成立。设 $\rho$ 把 $x_i$ 映到
$r_i \in R$ 的类。把 $x_i$ 替换为 $x_i-r_i$ 后，可以假设
$\rho(x_i) = 0$ 于 $R/\pi^4 R$ 中。这蕴含
$f_j(0) \in \pi^4R$，并且 $A \to \Lambda$ 把 $x_i$
映到某个 $\pi^4\lambda_i$，其中 $\lambda_i \in \Lambda$。写成
$$
f_j = f_j(0) + \sum\nolimits_{i = 1, \ldots, n} r_{ji} x_i + \text{高阶项}
$$
于是 $\partial f_j/\partial x_i$ 的常数项为 $r_{ji}$。
把 $\rho$ 应用于关于 $\pi$ 的 (\ref{smoothing-equation-strictly-standard-one})，可得
$$
\pi = \sum\nolimits_{I \subset \{1, \ldots, n\},\ |I| = c}
r_I \det(r_{ji})_{j = 1, \ldots, c,\ i \in I} \bmod \pi^4R
$$
其中某些 $r_I \in R$。因此有
$$
u\pi = \sum\nolimits_{I \subset \{1, \ldots, n\},\ |I| = c}
r_I \det(r_{ji})_{j = 1, \ldots, c,\ i \in I}
$$
其中某个 $u \in 1 + \pi^3R$。由《代数》
引理 \ref{algebra-lemma-matrix-left-inverse}，存在一个
$n \times c$ 矩阵 $(s_{ik})$，使得
$$
u\pi \delta_{jk} = \sum\nolimits_{i = 1, \ldots, n} r_{ji}s_{ik}\quad
\text{对所有 } j, k = 1, \ldots, c
$$
（Kronecker 符号）。引入辅助变量
$v_1, \ldots, v_c, w_1, \ldots, w_n$，并置
$$
h_i = x_i - \pi^2 \sum\nolimits_{j = 1, \ldots c} s_{ij} v_j - \pi^3 w_i
$$
下文将不再另行说明地使用等式
$$
R[x_1, \ldots, x_n, v_1, \ldots, v_c, w_1, \ldots, w_n]/
(h_1, \ldots, h_n) = R[v_1, \ldots, v_c, w_1, \ldots, w_n]
$$
。在
$R[x_1, \ldots, x_n, v_1, \ldots, v_c, w_1, \ldots, w_n]/
(h_1, \ldots, h_n)$ 中有
\begin{align*}
f_j & = f_j(x_1 - h_1, \ldots, x_n - h_n) \\
& =
\pi^2 \sum\nolimits_{k = 1}^c
\left(\sum\nolimits_{i = 1}^n r_{ji} s_{ik}\right) v_k
+
\pi^3 \sum\nolimits_{i = 1}^n r_{ji}w_i \bmod \pi^4 \\
& =
\pi^3 v_j + \pi^3 \sum\nolimits_{i = 1}^n r_{ji}w_i \bmod \pi^4
\end{align*}
，其中 $1 \leq j \leq c$。因此，可以选取元素
$g_j \in R[v_1, \ldots, v_c, w_1, \ldots, w_n]$，使得
$g_j = v_j + \sum r_{ji}w_i \bmod \pi$，并且在 $R$-代数
$R[x_1, \ldots, x_n, v_1, \ldots, v_c, w_1, \ldots, w_n]/
(h_1, \ldots, h_n)$ 中有 $f_j = \pi^3 g_j$。置
$$
B = R[x_1, \ldots, x_n, v_1, \ldots, v_c, w_1, \ldots, w_n]/
(f_1, \ldots, f_m, h_1, \ldots, h_n, g_1, \ldots, g_c).
$$
$A \to B$ 的同态是显然的。定义 $B \to \Lambda$，令
$x_i \mapsto \pi^4\lambda_i$、$v_i \mapsto 0$ 且
$w_i \mapsto \pi \lambda_i$。于是，元素 $f_j$ 和 $h_i$ 显然都映到
$\Lambda$ 中的零。此外，$g_i$ 映到 $\pi\Lambda$ 中的某个元素 $t$，
且 $\pi^3t = 0$（因为模掉由诸 $h$ 生成的理想后，
$f_i = \pi^3 g_i$）。所以，假设
$\text{Ann}_\Lambda(\pi) = \text{Ann}_\Lambda(\pi^2)$ 蕴含 $t=0$。
因而，只需证明关于光滑性的断言即可。

\medskip\noindent
注意，$B_\pi \cong A_\pi[v_1, \ldots, v_c]$，因为方程
$f_i=0$ 蕴含 $g_i=0$。又因假设 $\pi$ 在 $A$ 中相对于 $R$ 是严格标准元，
$A_\pi$ 在 $R$ 上光滑，见引理 \ref{smoothing-lemma-elkik}，故
$B_\pi$ 在 $R$ 上光滑。

\medskip\noindent
置 $B' = R[v_1, \ldots, v_c, w_1, \ldots, w_n]/(g_1, \ldots, g_c)$。
由于 $g_i = v_i + \sum r_{ji}w_i \bmod \pi$，可知
$B'/\pi B' = R/\pi R[w_1, \ldots, w_n]$。因此，由《代数》引理
\ref{algebra-lemma-localize-relative-complete-intersection} 和
\ref{algebra-lemma-flat-fibre-smooth}
，$R \to B'$ 在 $V(\pi)$ 的每一点都光滑且相对维数为 $n$
（第一个引理表明它在这些素理想处为 syntomic，因而特别是平坦的；
随后第二个引理表明它是光滑的）。

\medskip\noindent
设 $\mathfrak q \subset B$ 是一个满足 $\pi \in \mathfrak q$ 的素理想，
并且对某个 $r \in \mathfrak a$ 有 $r \not \in \mathfrak q$。
记 $\mathfrak q' = B' \cap \mathfrak q$。我们断言，满射 $B' \to B$
诱导局部环的同构 $(B')_{\mathfrak q'} \to B_\mathfrak q$；这将完成
引理的证明。注意，$B_\mathfrak q$ 是 $(B')_{\mathfrak q'}$ 对由
$f_{c + j}$（$j = 1, \ldots, m - c$）生成之理想的商。我们作两个观察：
第一，$f_{c+j}$ 在 $(B')_{\mathfrak q'}$ 中的像可被 $\pi^2$ 整除；
第二，由 (\ref{smoothing-equation-star})，$\pi f_{c+j}$ 在
$(B')_{\mathfrak q'}$ 中的像可以写成
$\sum b_{j_1 j_2} f_{c + j_1}f_{c + j_2}$。因此，每个
$\pi f_{c+j}$ 的像都属于由元素 $\pi^2 f_{c+j'}$ 生成的理想。
由于 $(B')_{\mathfrak q'}$ 是诺特局部环，故其中
$\pi f_{c+j}=0$，见《代数》
引理 \ref{algebra-lemma-intersect-powers-ideal-module-zero}。
又因 $R \to (B')_{\mathfrak q'}$ 平坦，可知
$$
\left(\text{Ann}_R(\pi^2)/\text{Ann}_R(\pi)\right)
\otimes_R (B')_{\mathfrak q'}
=
\text{Ann}_{(B')_{\mathfrak q'}}(\pi^2)/\text{Ann}_{(B')_{\mathfrak q'}}(\pi)
$$
由于 $r \in \mathfrak a$ 在 $(B')_{\mathfrak q'}$ 中可逆，
这个模必为零。因此，$f_{c+j}$ 在 $(B')_{\mathfrak q'}$ 中的像为零，
正合所需。
\end{proof}

\begin{lemma}
\label{smoothing-lemma-desingularize-strictly-standard}
设 $R$ 为诺特环，$\Lambda$ 为 $R$-代数。取 $\pi \in R$，并假设
$\text{Ann}_R(\pi) = \text{Ann}_R(\pi^2)$ 且
$\text{Ann}_\Lambda(\pi) = \text{Ann}_\Lambda(\pi^2)$。
设 $A \to \Lambda$ 和 $D \to \Lambda$ 为 $R$-代数同态，
其中 $A$ 和 $D$ 均为有限表示。假设
\begin{enumerate}
\item $\pi$ 在 $A$ 中相对于 $R$ 是严格标准元；
\item 存在与到 $\Lambda/\pi^4\Lambda$ 的同态相容的 $R$-代数同态
$A/\pi^4 A \to D/\pi^4 D$。
\end{enumerate}
则可找到 $R$-代数同态 $B \to \Lambda$，其中 $B$ 为有限表示，
以及与到 $\Lambda$ 的同态相容的 $R$-代数同态 $A \to B$ 和 $D \to B$，
使得 $H_{D/R}B \subset H_{B/D}$ 且 $H_{D/R}B \subset H_{B/R}$。
\end{lemma}

\begin{proof}
把引理 \ref{smoothing-lemma-desingularize} 应用于
$$
D \longrightarrow A \otimes_R D \longrightarrow \Lambda
$$
以及 $\pi$ 在 $D$ 中的像。由引理
\ref{smoothing-lemma-strictly-standard-base-change}，$\pi$ 在 $A \otimes_R D$
中相对于 $D$ 是严格标准元。取 (2) 中同态所诱导的同态作为截面
$\rho : (A \otimes_R D)/\pi^4 (A \otimes_R D) \to D/\pi^4 D$。
于是引理 \ref{smoothing-lemma-desingularize} 适用，并给出分解
$A \otimes_R D \to B \to \Lambda$，其中 $B$ 为有限表示，且
$\mathfrak a B \subset H_{B/D}$，这里
$$
\mathfrak a = \text{Ann}_D(\text{Ann}_D(\pi^2)/\text{Ann}_D(\pi)).
$$
对 $D$ 的任意素理想 $\mathfrak q$，若 $D_\mathfrak q$ 在 $R$ 上平坦，
则有
$\text{Ann}_{D_\mathfrak q}(\pi^2)/\text{Ann}_{D_\mathfrak q}(\pi) = 0$
这是因为元素的零化子与平坦基变换可交换，且我们假设了
$\text{Ann}_R(\pi) = \text{Ann}_R(\pi^2)$。由于 $D$ 是诺特环，
$\text{Ann}_D(\pi^2)/\text{Ann}_D(\pi)$ 是有限 $D$-模，
故取其零化子与局部化可交换。因此 $\mathfrak a \not \subset \mathfrak q$，
从而 $D \to B$ 在 $B$ 中位于 $\mathfrak q$ 之上的任意素理想处光滑。
凡使 $R \to D$ 光滑的 $D$ 中素理想 $\mathfrak q$，都有
$D_\mathfrak q$ 在 $R$ 上平坦，所以 $H_{D/R}B \subset H_{B/D}$。
最后的包含关系 $H_{D/R}B \subset H_{B/R}$ 源于光滑环同态的复合仍光滑
（《代数》，引理 \ref{algebra-lemma-compose-smooth}）。
\end{proof}

\begin{lemma}
\label{smoothing-lemma-desingularize-lifting-apply}
设 $R$ 为诺特环，$\Lambda$ 为 $R$-代数。取 $\pi \in R$，并假设
$\text{Ann}_R(\pi) = \text{Ann}_R(\pi^2)$ 且
$\text{Ann}_\Lambda(\pi) = \text{Ann}_\Lambda(\pi^2)$。
设 $A \to \Lambda$ 为 $R$-代数同态，其中 $A$ 为有限表示，
并假设 $\pi$ 在 $A$ 中相对于 $R$ 是严格标准元。设
$$
A/\pi^8A \to \bar C \to \Lambda/\pi^8\Lambda
$$
是一个分解，其中 $\bar C$ 为有限表示。则可以找到分解
$A \to B \to \Lambda$，其中 $B$ 为有限表示，
$R_\pi \to B_\pi$ 光滑，并且
$$
H_{\bar C/(R/\pi^8 R)} \cdot \Lambda/\pi^8\Lambda
\subset
\sqrt{H_{B/R} \Lambda} \bmod \pi^8\Lambda.
$$
\end{lemma}

\begin{proof}
应用引理 \ref{smoothing-lemma-lifting}，得到 $R \to D \to \Lambda$ 以及分解
$\bar C/\pi^4\bar C \to D/\pi^4 D \to \Lambda/\pi^4\Lambda$
，使 $R \to D$ 在任何不含 $\pi$ 的素理想处都光滑，
并且在任何位于 $\bar C/\pi^4\bar C$ 中使
$R/\pi^8 R \to \bar C$ 光滑之素理想上方的素理想处也光滑。
由引理 \ref{smoothing-lemma-desingularize-strictly-standard}，可以找到有限表示的
$R$-代数 $B$，以及分解 $A \to B \to \Lambda$ 和
$D \to B \to \Lambda$，使得 $H_{D/R}B \subset H_{B/R}$。
略去验证这确实给出了本引理所述问题的解。
\end{proof}












\section{预备：约化到基域}
\label{smoothing-section-reduction}

\noindent
本节应用前几节的引理，证明只需在基环为域时证明主结果，
见引理 \ref{smoothing-lemma-reduce-to-field}。

\begin{situation}
\label{smoothing-situation-global}
这里 $R \to \Lambda$ 是诺特环之间的正则环同态。
\end{situation}

\noindent
设 $R \to \Lambda$ 如情形 \ref{smoothing-situation-global} 所述。
若 $\Lambda$ 是光滑 $R$-代数的滤过余极限，则称
{\it $R \to \Lambda$ 满足 PT}。

\begin{lemma}
\label{smoothing-lemma-product}
设 $R_i \to \Lambda_i$（$i=1,2$）如情形
\ref{smoothing-situation-global} 所述。若 $R_i \to \Lambda_i$（$i=1,2$）
均满足 PT，则 $R_1 \times R_2 \to \Lambda_1 \times \Lambda_2$ 满足 PT。
\end{lemma}

\begin{proof}
略。提示：滤过余极限的乘积仍是滤过余极限。
\end{proof}

\begin{lemma}
\label{smoothing-lemma-delocalize-base}
设 $R \to A \to \Lambda$ 为环同态，其中 $A$ 在 $R$ 上为有限表示；
设 $S \subset R$ 为乘法集。给定分解
$S^{-1}A \to B' \to S^{-1}\Lambda$，其中 $B'$ 在 $S^{-1}R$ 上光滑。
则可找到分解 $A \to B \to \Lambda$，使某个 $s \in S$ 的像
在 $B$ 中相对于 $R$ 是初等标准元
（定义 \ref{smoothing-definition-strictly-standard}）。
\end{lemma}

\begin{proof}
先对 $S^{-1}R \to B'$ 应用引理
\ref{smoothing-lemma-smooth-standard-smooth}，从而可假设 $B'$ 在 $S^{-1}R$
上标准光滑。写
$A = R[x_1, \ldots, x_n]/(g_1, \ldots, g_t)$，并设 $x_i$
在 $\Lambda$ 中映到 $\lambda_i$。可以写成
$B' = S^{-1}R[x_1, \ldots, x_{n + m}]/(f_1, \ldots, f_c)$
，其中某个 $c \geq n$，使得
$\det(\partial f_j/\partial x_i)_{i, j = 1, \ldots, c}$
在 $B'$ 中可逆，并且 $A \to B'$ 由 $x_i \mapsto x_i$ 给出，
见引理 \ref{smoothing-lemma-standard-smooth-include-generators}。
将 $i>n$ 的 $x_i$ 乘以 $S$ 中某个元素，并相应修改方程 $f_j$ 后，
可假设对 $i>n$，$B' \to S^{-1}\Lambda$ 把 $x_i$ 映到某个
$\lambda_i/1$，其中 $\lambda_i \in \Lambda$。选取关系式
$$
1 =
a_0 \det(\partial f_j/\partial x_i)_{i, j = 1, \ldots, c}
+
\sum\nolimits_{j = 1, \ldots, c} a_jf_j
$$
，其中某些 $a_j \in S^{-1}R[x_1, \ldots, x_{n + m}]$。
由于 $S$ 的每个元素在 $B'$ 中都可逆，消去分母后，可以假设
$f_j, a_j \in R[x_1, \ldots, x_{n + m}]$，并且
$$
s_0 = a_0 \det(\partial f_j/\partial x_i)_{i, j = 1, \ldots, c}
+
\sum\nolimits_{j = 1, \ldots, c} a_jf_j
$$
，其中某个 $s_0 \in S$。由于 $g_j$ 在
$S^{-1}R[x_1, \ldots, x_{n + m}]/(f_1, \ldots, x_c)$
中的像为零，可以找到元素 $s_j \in S$，使得 $s_jg_j=0$ 于
$R[x_1, \ldots, x_{n + m}]/(f_1, \ldots, f_c)$.
由于 $f_j$ 在 $S^{-1}\Lambda$ 中的像为零，可以找到 $s'_j \in S$，
使 $s'_j f_j(\lambda_1, \ldots, \lambda_{n + m}) = 0$ 于
$\Lambda$ 中。考虑环
$$
B = R[x_1, \ldots, x_{n + m}]/
(s'_1f_1, \ldots, s'_cf_c, g_1, \ldots, g_t)
$$
以及分解 $A \to B \to \Lambda$，其中 $B \to \Lambda$ 由
$x_i \mapsto \lambda_i$ 给出。我们断言
$s = s_0s_1 \ldots s_ts'_1 \ldots s'_c$ 在 $B$ 中相对于 $R$
是初等标准元，这将完成证明。事实上，
$s_j g_j \in (f_1, \ldots, f_c)$，因而
$sg_j \in (s'_1f_1, \ldots, s'_cf_c)$。最后有
$$
a_0\det(\partial s'_jf_j/\partial x_i)_{i, j = 1, \ldots, c}
+
\sum\nolimits_{j = 1, \ldots, c}
(s'_1 \ldots \hat{s'_j} \ldots s'_c) a_j s'_jf_j
=
s_0s'_1\ldots s'_c
$$
，而它整除 $s$，正合所需。
\end{proof}

\begin{lemma}
\label{smoothing-lemma-reduce-to-field}
\begin{slogan}
证明 Popescu 逼近可约化到域上代数的情形
\end{slogan}
若在情形 \ref{smoothing-situation-global} 中每当 $R$ 是域时 PT 都成立，
则 PT 一般地成立。
\end{lemma}

\begin{proof}
假设在情形 \ref{smoothing-situation-global} 中，只要 $R$ 是域，PT 就成立。
任取情形 \ref{smoothing-situation-global} 中的 $R \to \Lambda$。注意，
$R/I \to \Lambda/I\Lambda$ 仍是诺特环之间的正则环同态，
见《代数详论》引理 \ref{more-algebra-lemma-regular-base-change}。
考虑理想的集合
$$
\mathcal{I} = \{I \subset R \mid R/I \to \Lambda/I\Lambda
\text{不满足 PT}\}
$$
我们必须证明 $\mathcal{I}$ 为空。若它非空，则因 $R$ 是诺特环，
其中存在极大元。把 $R$ 替换为 $R/I$，把 $\Lambda$ 替换为
$\Lambda/I$，便得到这样一种情形：对 $R$ 的任意非零理想 $I$，
$R/I \to \Lambda/I\Lambda$ 都满足 PT。特别地，应用命题
\ref{smoothing-proposition-lift} 可知 $R$ 是既约环。

\medskip\noindent
设 $A \to \Lambda$ 为 $R$-代数同态，其中 $A$ 为有限表示。
必须找到分解 $A \to B \to \Lambda$，使 $B$ 在 $R$ 上光滑，
见《代数》引理 \ref{algebra-lemma-when-colimit}。

\medskip\noindent
设 $S \subset R$ 为非零因子的集合，并考虑 $R$ 的全分式环
$Q=S^{-1}R$。已知 $Q = K_1 \times \ldots \times K_n$ 是若干域的乘积，
见《代数》引理 \ref{algebra-lemma-total-ring-fractions-no-embedded-points}
和 \ref{algebra-lemma-Noetherian-irreducible-components}。
由引理 \ref{smoothing-lemma-product} 及我们的假设，环同态
$S^{-1}R \to S^{-1}\Lambda$ 满足 PT。因此，可以找到分解
$S^{-1}A \to B' \to S^{-1}\Lambda$，其中 $B'$ 在 $S^{-1}R$ 上光滑。

\medskip\noindent
应用引理 \ref{smoothing-lemma-delocalize-base}，得到分解
$A \to B \to \Lambda$，使某个 $\pi \in S$ 在 $B$ 中相对于 $R$
是初等标准元。把 $A$ 替换为 $B$ 后，可假设 $\pi$ 在 $A$ 中是
初等标准元，因而是严格标准元。已知
$R/\pi^8R \to \Lambda/\pi^8\Lambda$ 满足 PT，所以可以找到分解
$R/\pi^8 R \to A/\pi^8A \to \bar C \to \Lambda/\pi^8\Lambda$
，其中 $R/\pi^8 R \to \bar C$ 光滑。由引理 \ref{smoothing-lemma-lifting}，
可以找到 $R$-代数同态 $D \to \Lambda$，其中 $D$ 在 $R$ 上光滑，
以及分解
$R/\pi^4 R \to A/\pi^4A \to D/\pi^4D \to \Lambda/\pi^4\Lambda$.
由引理 \ref{smoothing-lemma-desingularize-strictly-standard}，可以找到
$A \to B \to \Lambda$，其中 $B$ 在 $R$ 上光滑，从而完成证明。
\end{proof}










\section{局部技巧}
\label{smoothing-section-local}


\begin{situation}
\label{smoothing-situation-local}
给定诺特环 $R$、$R$-代数同态 $A \to \Lambda$，以及素理想
$\mathfrak q \subset \Lambda$。假设 $A$ 在 $R$ 上为有限表示。
在此情形下，记 $\mathfrak h_A = \sqrt{H_{A/R} \Lambda}$。
\end{situation}

\noindent
设 $R \to A \to \Lambda \supset \mathfrak q$ 如情形
\ref{smoothing-situation-local} 所述。若存在分解 $A \to B \to \Lambda$，
其中 $B$ 为有限表示，且
$\mathfrak h_A \subset \mathfrak h_B \not \subset \mathfrak q$，
则称 {\it $R \to A \to \Lambda \supset \mathfrak q$ 可消解}。
此时称分解 $A \to B \to \Lambda$ 为
{\it $R \to A \to \Lambda \supset \mathfrak q$ 的一个消解}。

\begin{lemma}
\label{smoothing-lemma-lift-solution}
设 $R \to A \to \Lambda \supset \mathfrak q$ 如情形
\ref{smoothing-situation-local} 所述。设 $r \geq 1$，并设
$\pi_1, \ldots, \pi_r \in R$ 的像属于 $\mathfrak q$。假设
\begin{enumerate}
\item 对 $i = 1, \ldots, r$ 有
$$
\text{Ann}_{R/(\pi_1^8, \ldots, \pi_{i - 1}^8)R}(\pi_i)
=
\text{Ann}_{R/(\pi_1^8, \ldots, \pi_{i - 1}^8)R}(\pi_i^2)
$$
且
$$
\text{Ann}_{\Lambda/(\pi_1^8, \ldots, \pi_{i - 1}^8)\Lambda}(\pi_i)
=
\text{Ann}_{\Lambda/(\pi_1^8, \ldots, \pi_{i - 1}^8)\Lambda}(\pi_i^2)
$$
\item 对 $i = 1, \ldots, r$，元素 $\pi_i$ 的像在 $A$ 中相对于 $R$
是严格标准元。
\end{enumerate}
那么，若
$$
R/(\pi_1^8, \ldots, \pi_r^8)R \to A/(\pi_1^8, \ldots, \pi_r^8)A
\to \Lambda/(\pi_1^8, \ldots, \pi_r^8)\Lambda \supset
\mathfrak q/(\pi_1^8, \ldots, \pi_r^8)\Lambda
$$
可消解，则 $R \to A \to \Lambda \supset \mathfrak q$ 也可消解。
\end{lemma}

\begin{proof}
对 $r$ 作归纳证明。

\medskip\noindent
先考虑 $r=1$。此时假设存在分解
$A/\pi_1^8 \to \bar C \to \Lambda/\pi_1^8$，它消解了模
$\pi_1^8$ 的情形。条件 (1) 和 (2) 正是应用引理
\ref{smoothing-lemma-desingularize-lifting-apply} 所需的假设。
因此，可以把消解 $\bar C$“提升”为
$R \to A \to \Lambda \supset \mathfrak q$ 的消解。

\medskip\noindent
再考虑 $r>1$。此时，把关于 $r-1$ 的归纳假设应用于情形
$R/\pi_1^8 \to A/\pi_1^8 \to \Lambda/\pi_1^8
\supset \mathfrak q/\pi_1^8\Lambda$.
注意，引理 \ref{smoothing-lemma-strictly-standard-base-change} 表明性质 (2)
在这里得到保持。
\end{proof}

\begin{lemma}
\label{smoothing-lemma-delocalize-weak}
\begin{reference}
\cite[Lemma 12.2]{swan} 或
\cite[Lemma 2]{popescu-GND}
\end{reference}
设 $R \to A \to \Lambda \supset \mathfrak q$ 如情形
\ref{smoothing-situation-local} 所述，并令 $\mathfrak p=R\cap\mathfrak q$。
假设 $\mathfrak q$ 是包含 $\mathfrak h_A$ 的极小素理想，并且
$R_\mathfrak p \to A_\mathfrak p \to \Lambda_\mathfrak q
\supset \mathfrak q\Lambda_\mathfrak q$ 可消解。
则存在分解 $A \to C \to \Lambda$，其中 $C$ 为有限表示，且
$H_{C/R} \Lambda \not \subset \mathfrak q$。
\end{lemma}

\begin{proof}
设 $A_\mathfrak p \to C \to \Lambda_\mathfrak q$ 是
$R_\mathfrak p \to A_\mathfrak p \to \Lambda_\mathfrak q
\supset \mathfrak q\Lambda_\mathfrak q$ 的一个消解。由
$\mathfrak q$ 是包含 $\mathfrak h_A$ 的极小素理想这一假设，
这意味着 $H_{C/R_\mathfrak p} \Lambda_\mathfrak q = \Lambda_\mathfrak q$。
由引理 \ref{smoothing-lemma-final-solve}，可假设 $C$ 在 $R_\mathfrak p$ 上光滑；
再由引理 \ref{smoothing-lemma-smooth-standard-smooth}，可假设 $C$ 在
$R_\mathfrak p$ 上标准光滑。写
$A = R[x_1, \ldots, x_n]/(g_1, \ldots, g_t)$，并设
$A \to \Lambda$ 由 $x_i \mapsto \lambda_i$ 给出。写
$C = R_\mathfrak p[x_1, \ldots, x_{n + m}]/(f_1, \ldots, f_c)$，
其中某个 $c \geq n$，使 $A \to C$ 把 $x_i$ 映到 $x_i$，并且
$\det(\partial f_j/\partial x_i)_{i, j = 1, \ldots, c}$
在 $C$ 中可逆，见引理 \ref{smoothing-lemma-standard-smooth-include-generators}。
消去分母后，可假设 $f_1,\ldots,f_c$ 属于
$R[x_1, \ldots, x_{n + m}]$。当然，
$\det(\partial f_j/\partial x_i)_{i, j = 1, \ldots, c}$
在 $R[x_1, \ldots, x_{n + m}]/(f_1, \ldots, f_c)$ 中未必可逆，
但将某个 $s_0 \in R$（$s_0 \not \in \mathfrak p$）取逆后便可逆。
由于 $g_j$ 经 $R[x_1, \ldots, x_n] \to A \to C$ 映到零，
可以找到 $s_j \in R$，其中 $s_j \not \in \mathfrak p$，使
$s_jg_j$ 在 $R[x_1, \ldots, x_{n + m}]/(f_1, \ldots, f_c)$ 中为零。
将 $f_j$ 写成 $f_j = F_j(x_1, \ldots, x_{n + m}, 1)$，其中多项式
$F_j \in R[x_1, \ldots, x_n, X_{n + 1}, \ldots, X_{n + m + 1}]$
关于 $X_{n + 1}, \ldots, X_{n + m + 1}$ 齐次。选取
$\lambda_{n + i} \in \Lambda$（$i = 1, \ldots, m + 1$），其中
$\lambda_{n + m + 1} \not \in \mathfrak q$，使 $x_{n+i}$ 在
$\Lambda_\mathfrak q$ 中映到
$\lambda_{n + i}/\lambda_{n + m + 1}$。于是
\begin{align*}
F_j(\lambda_1, \ldots, \lambda_{n + m + 1})
& =
(\lambda_{n + m + 1})^{\deg(F_j)} F_j(\lambda_1, \ldots, \lambda_n,
\frac{\lambda_{n + 1}}{\lambda_{n + m + 1}}, \ldots,
\frac{\lambda_{n + m}}{\lambda_{n + m + 1}}, 1) \\
& =
(\lambda_{n + m + 1})^{\deg(F_j)} f_j(\lambda_1, \ldots, \lambda_n,
\frac{\lambda_{n + 1}}{\lambda_{n + m + 1}}, \ldots,
\frac{\lambda_{n + m}}{\lambda_{n + m + 1}}) \\
& = 0
\end{align*}
在 $\Lambda_\mathfrak q$ 中成立。因此，可以找到
$\lambda_0 \in \Lambda$，其中 $\lambda_0 \not \in \mathfrak q$，使得
$\lambda_0 F_j(\lambda_1, \ldots, \lambda_{n + m + 1}) = 0$
在 $\Lambda$ 中成立。现在置 $B$ 为
$$
R[x_0, \ldots, x_{n + m + 1}]/
(g_1, \ldots, g_t, x_0F_1(x_1, \ldots, x_{n + m + 1}), \ldots,
x_0F_c(x_1, \ldots, x_{n + m + 1}))
$$
，并通过 $x_i \mapsto \lambda_i$ 将其映到 $\Lambda$。
令 $b$ 为 $x_0x_{n+m+1}s_0s_1\ldots s_t$ 在 $B$ 中的像。
则 $B_b$ 同构于
$$
R_{s_0s_1 \ldots s_t}[x_0, x_1, \ldots, x_{n + m + 1}, 1/x_0x_{n + m + 1}]/
(f_1, \ldots, f_c)
$$
，后者依构造在 $R$ 上光滑。由于 $b$ 的像不属于 $\mathfrak q$，
结论得证。
\end{proof}

\begin{lemma}
\label{smoothing-lemma-delocalize-height-zero}
设 $R \to A \to \Lambda \supset \mathfrak q$ 如情形
\ref{smoothing-situation-local} 所述，并令 $\mathfrak p=R\cap\mathfrak q$。假设
\begin{enumerate}
\item $\mathfrak q$ 是包含 $\mathfrak h_A$ 的极小素理想；
\item $R_\mathfrak p \to A_\mathfrak p \to \Lambda_\mathfrak q
\supset \mathfrak q\Lambda_\mathfrak q$ 可消解；
\item $\dim(\Lambda_\mathfrak q) = 0$.
\end{enumerate}
则 $R \to A \to \Lambda \supset \mathfrak q$ 可消解。
\end{lemma}

\begin{proof}
由 (3)，环 $\Lambda_\mathfrak q$ 是 Artin 局部环，因而
$\mathfrak q\Lambda_\mathfrak q$ 幂零。故对某个 $N>0$，有
$(\mathfrak h_A)^N \Lambda_\mathfrak q = 0$。于是存在
$\lambda \in \Lambda$，其中 $\lambda \not \in \mathfrak q$，使
$\lambda(\mathfrak h_A)^N=0$ 于 $\Lambda$ 中。设
$H_{A/R}=(a_1,\ldots,a_r)$，于是 $\lambda a_i^N=0$ 于 $\Lambda$ 中。
由引理 \ref{smoothing-lemma-delocalize-weak}，可以找到分解
$A \to C \to \Lambda$，其中 $C$ 为有限表示，且
$\mathfrak h_C \not \subset \mathfrak q$。写
$C=A[x_1,\ldots,x_n]/(f_1,\ldots,f_m)$。置
$$
B = A[x_1, \ldots, x_n, y_1, \ldots, y_r, z, t_{ij}]/
(f_j - \sum y_i t_{ij}, zy_i)
$$
其中 $t_{ij}$ 是 $rm$ 个变量。令 $t_{ij}$ 全部等于零，便得到同态
$B \to C[y_i,z]/(y_iz)$。同态 $B \to \Lambda$ 是复合
$B \to C[y_i,z]/(y_iz) \to \Lambda$，其中
$C[y_i,z]/(y_iz) \to \Lambda$ 延拓给定同态 $C \to \Lambda$，
把 $z$ 映到 $\lambda$，并把 $y_i$ 映到 $a_i^N$ 在 $\Lambda$ 中的像。

\medskip\noindent
我们断言，$B$ 是 $R \to A \to \Lambda \supset \mathfrak q$ 的一个解。
首先，$B_z$ 同构于 $C[y_1, \ldots, y_r,z,z^{-1}]$，因而是光滑的。
另一方面，
$B_{y_\ell} \cong A[x_i, y_i, y_\ell^{-1}, t_{ij}, i \not = \ell]$
在 $A$ 上光滑。因此，$z$ 以及各个 $a_\ell y_\ell$
（光滑同态的复合仍光滑）全都属于 $H_{B/R}$。引理得证。
\end{proof}




\section{可分剩余域}
\label{smoothing-section-separable}

\noindent
本节说明当剩余域扩张可分时如何求解一个局部问题。

\begin{lemma}[Ogoma]
\label{smoothing-lemma-ogoma}
设 $A$ 为诺特环，$M$ 为有限 $A$-模，$S\subset A$ 为乘法集。
若 $\pi\in A$ 且
$\Ker(\pi : S^{-1}M \to S^{-1}M) =
\Ker(\pi^2 : S^{-1}M \to S^{-1}M)$
，则存在 $s\in S$，使得对任意 $n>0$ 都有
$\Ker(s^n\pi : M \to M) = \Ker((s^n\pi)^2 : M \to M)$.
\end{lemma}

\begin{proof}
令 $K=\Ker(\pi:M\to M)$、
$K'=\{m\in M\mid \pi^2m=0\text{ 于 }S^{-1}M\text{ 中}\}$，
并令 $Q=K'/K$。由假设，$S^{-1}Q=0$。由于 $A$ 是诺特环，
$Q$ 是有限 $A$-模。因此可以找到 $s\in S$ 使 $s$ 零化 $Q$；
这个 $s$ 即合要求。
\end{proof}

\begin{lemma}
\label{smoothing-lemma-find-sequence}
设 $\Lambda$ 为诺特环，$I\subset\Lambda$ 为理想，
且 $I\subset\mathfrak q$，其中 $\mathfrak q$ 为素理想。
设 $n,e$ 为正整数。假设
$\mathfrak q^n\Lambda_\mathfrak q\subset I\Lambda_\mathfrak q$，
且 $\Lambda_\mathfrak q$ 是维数为 $d$ 的正则局部环。
则存在 $n>0$ 以及 $\pi_1,\ldots,\pi_d\in\Lambda$，使得
\begin{enumerate}
\item $(\pi_1, \ldots, \pi_d)\Lambda_\mathfrak q =
\mathfrak q\Lambda_\mathfrak q$；
\item $\pi_1^n, \ldots, \pi_d^n \in I$；
\item 对 $i=1,\ldots,d$ 有
$$
\text{Ann}_{\Lambda/(\pi_1^e, \ldots, \pi_{i - 1}^e)\Lambda}(\pi_i) =
\text{Ann}_{\Lambda/(\pi_1^e, \ldots, \pi_{i - 1}^e)\Lambda}(\pi_i^2).
$$
\end{enumerate}
\end{lemma}

\begin{proof}
置 $S=\Lambda\setminus\mathfrak q$，于是
$\Lambda_\mathfrak q=S^{-1}\Lambda$。先选取满足 (1) 的
$\pi_1,\ldots,\pi_d$；这因 $\Lambda_\mathfrak q$ 正则而可行。
由假设，$\pi_i^n\in I\Lambda_\mathfrak q$，故可找到
$s_1,\ldots,s_d\in S$，使 $s_i\pi_i^n\in I$。以
$s_i\pi_i$ 替换 $\pi_i$ 后便得到 (2)。注意，继续乘以 $S$ 中元素
仍保持 (1) 和 (2)。假设对某个 $t\in\{0,\ldots,d\}$，
(3) 对 $i=1,\ldots,t$ 成立。注意，
$\pi_1,\ldots,\pi_d$ 是 $S^{-1}\Lambda$ 中的正则序列，
见《代数》引理 \ref{algebra-lemma-regular-ring-CM}。
特别地，由《代数》引理 \ref{algebra-lemma-regular-sequence-powers}，
$\pi_1^e,\ldots,\pi_t^e,\pi_{t+1}$ 是
$S^{-1}\Lambda=\Lambda_\mathfrak q$ 中的正则序列。因此
$$
\text{Ann}_{S^{-1}\Lambda/(\pi_1^e, \ldots, \pi_{i - 1}^e)}(\pi_i) =
\text{Ann}_{S^{-1}\Lambda/(\pi_1^e, \ldots, \pi_{i - 1}^e)}(\pi_i^2).
$$
由引理 \ref{smoothing-lemma-ogoma}，取某个 $s\in S$，以
$s\pi_{t+1}$ 替换 $\pi_{t+1}$ 后，便可使 (3) 对 $i=t+1$ 成立。
对 $t$ 归纳，即得到满足 (1)、(2) 和 (3) 的序列。
\end{proof}

\begin{lemma}
\label{smoothing-lemma-resolve-special}
设 $k \to A \to \Lambda \supset \mathfrak q$ 如情形
\ref{smoothing-situation-local} 所述，其中
\begin{enumerate}
\item $k$ 是域；
\item $\Lambda$ 是诺特环；
\item $\mathfrak q$ 是包含 $\mathfrak h_A$ 的极小素理想；
\item $\Lambda_\mathfrak q$ 是正则局部环；
\item 域扩张 $\kappa(\mathfrak q)/k$ 可分。
\end{enumerate}
则 $k \to A \to \Lambda \supset \mathfrak q$ 可消解。
\end{lemma}

\begin{proof}
置 $d=\dim\Lambda_\mathfrak q$ 及 $R=k[x_1,\ldots,x_d]$。
选取 $n>0$，使
$\mathfrak q^n\Lambda_\mathfrak q\subset\mathfrak h_A\Lambda_\mathfrak q$；
这因 $\mathfrak q$ 是包含 $\mathfrak h_A$ 的极小素理想而可行。
选取 $H_{A/R}$ 的生成元 $a_1,\ldots,a_r$。置
$$
B = A[x_1, \ldots, x_d, z_{ij}]/(x_i^n - \sum z_{ij}a_j)
$$
每个 $B_{a_j}$ 都在 $R$ 上光滑：它是
$A_{a_j}[x_1,\ldots,x_d]$ 上的多项式代数，而 $A_{a_j}$ 在 $k$ 上光滑。
因此 $B_{x_i}$ 在 $R$ 上光滑。令 $B\to C$ 为引理
\ref{smoothing-lemma-improve-presentation} 中构造的 $R$-代数同态；该构造还给出
$R$-代数缩回 $C\to B$，特别给出与上述图表相容的同态
$C\to\Lambda$。依构造，$C_{x_i}$ 是光滑 $R$-代数，且
$\Omega_{C_{x_i}/R}$ 自由。因此，可以找到 $c>0$，使 $x_i^c$
在 $C/R$ 中是严格标准元，见引理 \ref{smoothing-lemma-compare-standard}。
现在按照引理 \ref{smoothing-lemma-find-sequence} 选取
$\pi_1,\ldots,\pi_d\in\Lambda$，
其中 $n=n$、$e=8c$、$\mathfrak q=\mathfrak q$ 且
$I=\mathfrak h_A$。写
$\pi_i^n=\sum\lambda_{ij}a_j$，其中某些 $\pi_{ij}\in\Lambda$。
存在同态 $B\to\Lambda$，由 $x_i\mapsto\pi_i$ 及
$z_{ij}\mapsto\lambda_{ij}$ 给出。置 $R=k[x_1,\ldots,x_d]$。图表为
$$
\xymatrix{
R \ar[r] & B \ar[rd] \\
k \ar[u] \ar[r] & A \ar[u] \ar[r] & \Lambda
}
$$
现在把引理 \ref{smoothing-lemma-lift-solution} 应用于
$R \to C \to \Lambda \supset \mathfrak q$ 以及 $R$ 中的元素序列
$x_1^c,\ldots,x_d^c$。假设 (2) 显然成立。对 $R$ 而言，
直接检验可知假设 (1) 成立；对 $\Lambda$ 而言，它由我们对
$\pi_1,\ldots,\pi_d$ 的选择成立。（注意，若
$\text{Ann}_\Lambda(\pi)=\text{Ann}_\Lambda(\pi^2)$，则对所有
$c>0$ 都有 $\text{Ann}_\Lambda(\pi)=\text{Ann}_\Lambda(\pi^c)$。）
因此，只需消解
$$
R/(x_1^e, \ldots, x_d^e) \to C/(x_1^e, \ldots, x_d^e) \to
\Lambda/(\pi_1^e, \ldots, \pi_d^e) \supset
\mathfrak q/(\pi_1^e, \ldots, \pi_d^e)
$$
，其中 $e=8c$。由引理 \ref{smoothing-lemma-delocalize-height-zero}，
只需在 $\mathfrak q$ 处局部化后消解它。由于
$x_1,\ldots,x_d$ 在 $\Lambda_\mathfrak q$ 中映成正则序列，
$R_\mathfrak p\to\Lambda_\mathfrak q$ 平坦，见《代数》
引理 \ref{algebra-lemma-flat-over-regular}。所以
$$
R_\mathfrak p/(x_1^e, \ldots, x_d^e) \to
\Lambda_\mathfrak q/(\pi_1^e, \ldots, \pi_d^e)
$$
是 Artin 局部环之间的平坦环同态。此外，由假设，该同态在剩余域上
诱导可分域扩张。因此，由《代数》引理
\ref{algebra-lemma-colimit-syntomic} 和命题 \ref{smoothing-proposition-lift}，
这个同态是光滑代数的滤过余极限。所需解的存在性由《代数》
引理 \ref{algebra-lemma-when-colimit} 得出。
\end{proof}






\section{不可分剩余域}
\label{smoothing-section-inseparable}

\noindent
本节说明当剩余域扩张不可分时如何求解一个局部问题。

\begin{lemma}
\label{smoothing-lemma-helper}
设 $k$ 为特征 $p>0$ 的域，
$(\Lambda,\mathfrak m,K)$ 为 Artin 局部 $k$-代数。
假设 $\dim H_1(L_{K/k})<\infty$。则 $\Lambda$ 是 Artin 局部
$k$-代数 $A$ 的滤过余极限，其中每个同态 $A\to\Lambda$ 都平坦，
$\mathfrak m_A\Lambda=\mathfrak m$，且 $A$ 在 $k$ 上本质有限型。
\end{lemma}

\begin{proof}
注意，$A\to\Lambda$ 的平坦性蕴含该同态为单射，故本引理实际说明
$\Lambda$ 是这类子环 $A\subset\Lambda$ 的有向并。令 $n$ 为满足
$\mathfrak m^n=0$ 的最小整数。对 $n$ 归纳证明本引理。
$n=1$ 的情形显然，因为任一域扩张都是有限生成域扩张的并。

\medskip\noindent
选取生成 $\mathfrak m$ 的元素
$\lambda_1,\ldots,\lambda_d\in\mathfrak m$。由于 $K$ 在
$\mathbf{F}_p$ 上形式光滑（见《代数》引理
\ref{algebra-lemma-formally-smooth-extensions-easy}），可以找到环同态
$\sigma:K\to\Lambda$，它是商同态 $\Lambda\to K$ 的截面。
一般而言，$\sigma$ {\bf 不是} $k$-代数同态。给定 $\sigma$，定义
$$
\Psi_\sigma : K[x_1, \ldots, x_d] \longrightarrow \Lambda
$$
：在 $K$ 的元素上使用 $\sigma$，并令 $x_i\mapsto\lambda_i$。
断言：存在 $\sigma:K\to\Lambda$ 以及在 $k$ 上有限生成的子域
$k\subset F\subset K$，使 $k$ 在 $\Lambda$ 中的像包含于
$\Psi_\sigma(F[x_1,\ldots,x_d])$。

\medskip\noindent
对满足 $\mathfrak m^n=0$ 的最小整数 $n$ 作归纳，以证明该断言。
$n=1$ 时显然。若 $n>1$，置 $I=\mathfrak m^{n-1}$ 及
$\Lambda'=\Lambda/I$。由归纳假设，可认为已给定
$\sigma':K\to\Lambda'$ 以及有限生成的 $k\subset F'\subset K$，
使 $k\to\Lambda\to\Lambda'$ 的像包含于
$A'=\Psi_{\sigma'}(F'[x_1,\ldots,x_d])$。记
$\tau':k\to A'$ 为所诱导的同态。选取 $\sigma'$ 的一个提升
$\sigma:K\to\Lambda$；这可由前述 $K/\mathbf{F}_p$ 的形式光滑性做到。
留意，稍后可用某个导子 $D:K\to I$ 把 $\sigma$ 改为 $\sigma+D$。
置 $A=F[x_1,\ldots,x_d]/(x_1,\ldots,x_d)^n$。于是
$\Psi_\sigma$ 诱导环同态 $\Psi_\sigma:A\to\Lambda$；它与商同态
$\Lambda\to\Lambda'$ 的复合诱导一个核为幂零理想的满射 $A\to A'$。
选取 $\tau'$ 的提升 $\tau:k\to A$；这可由
$k/\mathbf{F}_p$ 的形式光滑性做到。于是得到两个从 $k$ 到
$\Lambda$ 的同态，即 $\Psi_\sigma\circ\tau:k\to\Lambda$ 和给定同态
$i:k\to\Lambda$。它们模 $I$ 相同，故其差为导子
$\theta=i-\Psi_\sigma\circ\tau:k\to I$。注意，若把 $\sigma$
改为 $\sigma+D$，则 $\theta$ 相应改为 $\theta-D|_k$。

\medskip\noindent
在 $k$ 中选取一族元素 $\{y_j\}_{j\in J}$，使其微分
$\text{d}y_j$ 构成 $\Omega_{k/\mathbf{F}_p}$ 的一组基。
对 $\mathbf{F}_p\subset k\subset K$，Jacobi--Zariski 序列为
$$
0 \to H_1(L_{K/k}) \to \Omega_{k/\mathbf{F}_p} \otimes K \to
\Omega_{K/\mathbf{F}_p} \to \Omega_{K/k} \to 0
$$
由于 $\dim H_1(L_{K/k})<\infty$，可以找到有限子集
$J_0\subset J$，使第一个同态的像包含于
$\bigoplus_{j\in J_0}K\text{d}y_j$。因此，元素
$\text{d}y_j$（$j\in J\setminus J_0$）在
$\Omega_{K/\mathbf{F}_p}$ 中的像在 $K$ 上线性无关。故可选取
$D:K\to I$，使 $\theta-D|_k=\xi\circ\text{d}$，其中 $\xi$ 是复合
$$
\Omega_{k/\mathbf{F}_p} = \bigoplus\nolimits_{j \in J} k \text{d}y_j
\longrightarrow \bigoplus\nolimits_{j \in J_0} k \text{d}y_j
\longrightarrow I
$$
对 $j\in J_0$，令 $f_j=\xi(\text{d}y_j)\in I$。按上面所述，
把 $\sigma$ 改为 $\sigma+D$。于是，对 $a\in k$，若在
$\Omega_{k/\mathbf{F}_p}$ 中有
$\text{d}a=\sum a_j\text{d}y_j$，则
$\theta(a)=\sum_{j\in J_0}a_jf_j$。注意，$I$ 由
$\lambda_1,\ldots,\lambda_d$ 中总次数
$|E|=\sum e_i=n-1$ 的单项式
$\lambda^E=\lambda_1^{e_1}\ldots\lambda_d^{e_d}$ 生成。写
$f_j=\sum_Ec_{j,E}\lambda^E$，其中 $c_{j,E}\in K$。
把 $F'$ 替换为 $F=F'(c_{j,E})$，断言即得证。

\medskip\noindent
按断言选取 $\sigma$ 和 $F$。$\Psi_\sigma$ 的核由有限多个多项式
$g_1,\ldots,g_t\in K[x_1,\ldots,x_d]$ 生成。向 $F$ 中添加有限多个
元素将其扩大后，可假设这些多项式的系数都在 $F$ 中。此时同态
$A = F[x_1, \ldots, x_d]/(g_1, \ldots, g_t) \to
K[x_1, \ldots, x_d]/(g_1, \ldots, g_t) = \Lambda$ 显然平坦。
由该断言，$A$ 是 $\Lambda$ 的 $k$-子代数。由于 $K$ 是这些子域
$F$ 的滤过并，$\Lambda$ 显然是这些代数的滤过余极限。最后，
由《代数》引理
\ref{algebra-lemma-essentially-of-finite-type-into-artinian-local}，
这些代数在 $k$ 上本质有限型。
\end{proof}

\begin{lemma}
\label{smoothing-lemma-solution-modulo}
设 $k$ 为特征 $p>0$ 的域，$\Lambda$ 为几何正则的诺特
$k$-代数，$\mathfrak q\subset\Lambda$ 为素理想。设 $n\geq1$ 为整数，
且 $E\subset\Lambda_\mathfrak q/\mathfrak q^n\Lambda_\mathfrak q$
为有限子集。则可以找到 $m\geq0$ 以及
$\varphi:k[y_1,\ldots,y_m]\to\Lambda$，使得
\begin{enumerate}
\item 令 $\mathfrak p=\varphi^{-1}(\mathfrak q)$，则
$\mathfrak q\Lambda_\mathfrak q=\mathfrak p\Lambda_\mathfrak q$，且
$k[y_1,\ldots,y_m]_\mathfrak p\to\Lambda_\mathfrak q$ 平坦；
\item 存在由 Artin 局部环同态组成的分解
$$
k[y_1, \ldots, y_m]_\mathfrak p/\mathfrak p^n k[y_1, \ldots, y_m]_\mathfrak p
\to D \to
\Lambda_\mathfrak q/\mathfrak q^n\Lambda_\mathfrak q
$$
，其中第一个箭头本质光滑，第二个箭头平坦；
\item 模 $\mathfrak q^n\Lambda_\mathfrak q$ 意义下，$E$ 包含于 $D$。
\end{enumerate}
\end{lemma}

\begin{proof}
置 $\bar\Lambda=\Lambda_\mathfrak q/\mathfrak q^n\Lambda_\mathfrak q$。
由《代数详论》命题
\ref{more-algebra-proposition-characterization-geometrically-regular}.
，有 $\dim H_1(L_{\kappa(\mathfrak q)/k})<\infty$。
按照引理 \ref{smoothing-lemma-helper}，选取含 $E$ 的 $A\subset\bar\Lambda$，
使 $A$ 是在 $k$ 上本质有限型的 Artin 局部环，$A\to\bar\Lambda$
平坦，且 $\mathfrak m_A$ 生成 $\bar\Lambda$ 的极大理想。
记剩余域 $F=A/\mathfrak m_A$，于是 $k\subset F\subset K$。
选取 $\lambda_1,\ldots,\lambda_t\in\Lambda$，它们在
$\bar\Lambda$ 中映到 $A$ 的元素，并且
$\text{d}\lambda_1,\ldots,\text{d}\lambda_t$ 的像构成
$\Omega_{F/k}$ 的一组基。考虑把 $y_j$ 映到 $\lambda_j$ 的同态
$\varphi':k[y_1,\ldots,y_t]\to\Lambda$，并置
$\mathfrak p'=(\varphi')^{-1}(\mathfrak q)$。由《代数详论》引理
\ref{more-algebra-lemma-geometrically-regular-over-field}
，环同态 $k[y_1,\ldots,y_t]_{\mathfrak p'}\to\Lambda_\mathfrak q$
平坦，且 $\Lambda_\mathfrak q/\mathfrak p'\Lambda_\mathfrak q$ 正则。
因此，可以再选取元素 $\lambda_{t+1},\ldots,\lambda_m\in\Lambda$；
它们在 $\bar\Lambda$ 中映入 $A$，且在
$\Lambda_\mathfrak q/\mathfrak p'\Lambda_\mathfrak q$ 中映成一组
正则参数系。由此得到满足性质 (1) 的
$\varphi:k[y_1,\ldots,y_m]\to\Lambda$，并且
$k[y_1, \ldots, y_m]_\mathfrak p/\mathfrak p^n k[y_1, \ldots, y_m]_\mathfrak p
\to \bar\Lambda$
经由 $A$ 分解。因此，
$k[y_1, \ldots, y_m]_\mathfrak p/\mathfrak p^n k[y_1, \ldots, y_m]_\mathfrak p
\to A$ 由《代数》引理
\ref{algebra-lemma-flatness-descends-more-general} 可知是平坦的。
依构造，剩余域扩张 $F/\kappa(\mathfrak p)$ 有限生成，且
$\Omega_{F/\kappa(\mathfrak p)}=0$。故由《代数详论》引理
\ref{more-algebra-lemma-cartier-equality}，它是有限可分扩张。因此
$k[y_1, \ldots, y_m]_\mathfrak p/\mathfrak p^n k[y_1, \ldots, y_m]_\mathfrak p
\to A$
由《代数》引理
\ref{algebra-lemma-essentially-of-finite-type-into-artinian-local}
可知是有限的。最后，由《代数》引理
\ref{algebra-lemma-characterize-etale}，它是平展的。
平展环同态当然本质光滑，故结论得证。
\end{proof}

\begin{lemma}
\label{smoothing-lemma-enlarge-solution-modulo}
设 $\varphi:k[y_1,\ldots,y_m]\to\Lambda$、$n$、$\mathfrak q$、
$\mathfrak p$ 以及
$$
k[y_1, \ldots, y_m]_\mathfrak p/\mathfrak p^n \to
D \to \Lambda_\mathfrak q/\mathfrak q^n \Lambda_\mathfrak q
$$
如引理 \ref{smoothing-lemma-solution-modulo} 所述。则对任意
$\lambda\in\Lambda\setminus\mathfrak q$，存在整数 $q>0$ 以及分解
$$
k[y_1, \ldots, y_m]_\mathfrak p/\mathfrak p^n \to
D \to D' \to \Lambda_\mathfrak q/\mathfrak q^n \Lambda_\mathfrak q
$$
，使 $D\to D'$ 是 Artin 局部环之间的本质光滑同态，
最后一个箭头平坦，并且 $\lambda^q$ 属于 $D'$。
\end{lemma}

\begin{proof}
置 $\bar\Lambda=\Lambda_\mathfrak q/\mathfrak q^n\Lambda_\mathfrak q$，
并令 $\bar\lambda$ 为 $\lambda$ 在 $\bar\Lambda$ 中的像，
$\alpha\in\kappa(\mathfrak q)$ 为 $\lambda$ 在剩余域中的像。
设 $k\subset F\subset\kappa(\mathfrak q)$ 为 $D$ 的剩余域。
若 $\alpha\in F$，则可找到 $x\in D$，使
$x\bar\lambda=1\bmod\mathfrak q$。若 $q$ 可被 $p$ 整除，则
$(x\bar\lambda)^q=1\bmod(\mathfrak q)^q$，故
$\bar\lambda^q\in D$。若 $\alpha$ 在 $F$ 上超越，则可取
$D'=(D[\bar\lambda])_\mathfrak m$，即把由 $D$ 和 $\bar\lambda$
生成的子环在
$\mathfrak m=D[\bar\lambda]\cap\mathfrak q\bar\Lambda$ 处局部化。
这是可行的，因为此时 $D[\bar\lambda]$ 实际上是 $D$ 上的多项式代数。
最后，若 $\lambda\bmod\mathfrak q$ 在 $F$ 上代数，则可以找到
$p$ 的幂 $q$，使 $\alpha^q$ 在 $F$ 上可分代数，
见《域》一章第 \ref{fields-section-algebraic} 节。
注意，$D$ 和 $\bar\Lambda$ 都是 Hensel 局部环，见《代数》引理
\ref{algebra-lemma-local-dimension-zero-henselian}。令 $D\to D'$ 为
有限平展扩张，其剩余域扩张为 $F(\alpha^q)/F$，见《代数》引理
\ref{algebra-lemma-henselian-cat-finite-etale}。由于 $\bar\Lambda$
是 Hensel 环，且其剩余域包含 $F(\alpha^q)$，可以找到分解
$D'\to\bar\Lambda$。由论证的第一部分可知，对某个 $q'>0$，有
$\bar\lambda^{qq'}\in D'$。
\end{proof}

\begin{lemma}
\label{smoothing-lemma-resolve-general}
设 $k \to A \to \Lambda \supset \mathfrak q$ 如情形
\ref{smoothing-situation-local} 所述，其中
\begin{enumerate}
\item $k$ 是特征 $p>0$ 的域；
\item $\Lambda$ 是诺特环，且在 $k$ 上几何正则；
\item $\mathfrak q$ 是包含 $\mathfrak h_A$ 的极小素理想。
\end{enumerate}
则 $k \to A \to \Lambda \supset \mathfrak q$ 可消解。
\end{lemma}

\begin{proof}
按给定顺序完成以下步骤即可证明本引理。下面将逐一说明这些步骤。
\begin{enumerate}
\item
\label{smoothing-item-power}
选取整数 $N>0$，使
$\mathfrak q^N\Lambda_\mathfrak q \subset H_{A/k}\Lambda_\mathfrak q$.
\item
\label{smoothing-item-generators}
选取理想 $H_{A/R}$ 的生成元 $a_1,\ldots,a_t\in A$。
\item
\label{smoothing-item-dimension}
置 $d=\dim(\Lambda_\mathfrak q)$。
\item
\label{smoothing-item-standardizer}
置 $B=A[x_1,\ldots,x_d,z_{ij}]/(x_i^{2N}-\sum z_{ij}a_j)$。
\item
\label{smoothing-item-elkik}
把 $B$ 看作 $k[x_1,\ldots,x_d]$-代数，并令 $B\to C$ 如引理
\ref{smoothing-lemma-improve-presentation} 所述。我们还得到截面 $C\to B$。
\item
\label{smoothing-item-strictly-standard}
选取 $c>0$，使每个 $x_i^c$ 在 $C$ 中相对于
$k[x_1,\ldots,x_d]$ 都是严格标准元。
\item
\label{smoothing-item-set-n}
置 $n=N+dc$ 及 $e=8c$。
\item
\label{smoothing-item-set-E}
令 $E\subset\Lambda_\mathfrak q/\mathfrak q^n\Lambda_\mathfrak q$
为 $A$ 作为 $k$-代数的一组生成元的像。
\item
\label{smoothing-item-NP}
选取整数 $m$、$k$-代数同态
$\varphi : k[y_1, \ldots, y_m] \to \Lambda$
以及由 Artin 局部环组成的分解
$$
k[y_1, \ldots, y_m]_\mathfrak p/\mathfrak p^n k[y_1, \ldots, y_m]_\mathfrak p
\to D \to
\Lambda_\mathfrak q/\mathfrak q^n\Lambda_\mathfrak q
$$
，使第一个箭头本质光滑，第二个箭头平坦，$E$ 包含于 $D$；并且令
$\mathfrak p=\varphi^{-1}(\mathfrak q)$ 时，同态
$k[y_1,\ldots,y_m]_\mathfrak p\to\Lambda_\mathfrak q$ 平坦，且
$\mathfrak p\Lambda_\mathfrak q=\mathfrak q\Lambda_\mathfrak q$。
\item
\label{smoothing-item-choose-pii}
选取 $\pi_1,\ldots,\pi_d\in\mathfrak p$，它们映成
$k[y_1,\ldots,y_m]_\mathfrak p$ 的一组正则参数系。
\item
\label{smoothing-item-set-R}
令 $R=k[y_1,\ldots,y_m,t_1,\ldots,t_m]$ 且 $\gamma_i=\pi_it_i$。
\item
\label{smoothing-item-modify-pii}
必要时修改 $\pi_i$ 的选择，使得对 $i=1,\ldots,d$ 有
$$
\text{Ann}_{R/(\gamma_1^e, \ldots, \gamma_{i - 1}^e)R}(\gamma_i)
=
\text{Ann}_{R/(\gamma_1^e, \ldots, \gamma_{i - 1}^e)R}(\gamma_i^2)
$$
\item
\label{smoothing-item-choose-deltai}
存在 $\delta_1,\ldots,\delta_d\in\Lambda$，其中
$\delta_i \not \in \mathfrak q$，以及分解
$D \to D' \to \Lambda_\mathfrak q/\mathfrak q^n\Lambda_\mathfrak q$
，其中 $D'$ 是 Artin 局部环，$D\to D'$ 本质光滑，且同态
$D'\to\Lambda_\mathfrak q/\mathfrak q^n\Lambda_\mathfrak q$ 平坦；
令 $\pi_i'=\delta_i\pi_i$，则对 $i=1,\ldots,d$ 有
\begin{enumerate}
\item 在 $\Lambda$ 中，$(\pi_i')^{2N}=\sum a_j\lambda_{ij}$，其中
$\lambda_{ij}\bmod\mathfrak q^n\Lambda_\mathfrak q$ 是 $D'$ 的元素；
\item $\text{Ann}_{\Lambda/({\pi'}_1^e, \ldots, {\pi'}_{i - 1}^e)}({\pi'}_i) =
\text{Ann}_{\Lambda/({\pi'}_1^e, \ldots, {\pi'}_{i - 1}^e)}({\pi'}_i^2)$；
\item $\delta_i\bmod\mathfrak q^n\Lambda_\mathfrak q$ 是 $D'$ 的元素。
\end{enumerate}
\item
\label{smoothing-item-map-B-Lambda}
令 $x_i\mapsto\pi_i'$、$z_{ij}\mapsto\lambda_{ij}$，以此定义
$B\to\Lambda$。把 $B\to\Lambda$ 与缩回 $C\to B$ 复合，
定义 $C\to\Lambda$。
\item
\label{smoothing-item-set-map-R}
定义 $R\to\Lambda$：在 $k[y_1,\ldots,y_m]$ 上取 $\varphi$，
并令 $t_i\mapsto\delta_i$。再引入同态
$$
k[x_1, \ldots, x_d]
\longrightarrow
R = k[y_1, \ldots, y_m, t_1, \ldots, t_d]
$$
，令 $x_i\mapsto\gamma_i=\pi_it_i$。
\item
\label{smoothing-item-first-solve}
只需消解
$$
R
\to
C \otimes_{k[x_1, \ldots, x_d]} R
\to
\Lambda \supset \mathfrak q
$$
\item
\label{smoothing-item-set-I}
置 $I=(\gamma_1^e,\ldots,\gamma_d^e)\subset R$。
\item
\label{smoothing-item-second-resolve}
只需消解
$$
R/I
\to
C \otimes_{k[x_1, \ldots, x_d]} R/I
\to
\Lambda/I\Lambda \supset \mathfrak q/I\Lambda
$$
\item
\label{smoothing-item-set-r}
记 $\mathfrak q$ 在
$R=k[y_1,\ldots,y_m,t_1,\ldots,t_d]$ 中的逆像为
$\mathfrak r\subset R$。
\item
\label{smoothing-item-third-resolve}
只需消解
$$
(R/I)_\mathfrak r \to
C \otimes_{k[x_1, \ldots, x_d]} (R/I)_\mathfrak r \to
\Lambda_\mathfrak q/I\Lambda_\mathfrak q
\supset
\mathfrak q\Lambda_\mathfrak q/I\Lambda_\mathfrak q
$$
\item
\label{smoothing-item-fourth-resolve}
在 $k[y_1,\ldots,y_m]$ 中置 $J=(\pi_1^e,\ldots,\pi_d^e)$。
\item
\label{smoothing-item-fifth-resolve}
只需消解
$$
(R/JR)_\mathfrak p \to
C \otimes_{k[x_1, \ldots, x_d]} (R/JR)_\mathfrak p \to
\Lambda_\mathfrak q/J\Lambda_\mathfrak q
\supset
\mathfrak q\Lambda_\mathfrak q/J\Lambda_\mathfrak q
$$
\item
\label{smoothing-item-sixth-resolve}
只需消解
$$
(R/\mathfrak p^nR)_\mathfrak p \to
C \otimes_{k[x_1, \ldots, x_d]} (R/\mathfrak p^nR)_\mathfrak p \to
\Lambda_\mathfrak q/\mathfrak q^n\Lambda_\mathfrak q
\supset
\mathfrak q\Lambda_\mathfrak q/\mathfrak q^n\Lambda_\mathfrak q
$$
\item
\label{smoothing-item-seventh-resolve}
只需消解
$$
(R/\mathfrak p^nR)_\mathfrak p \to
B \otimes_{k[x_1, \ldots, x_d]} (R/\mathfrak p^nR)_\mathfrak p \to
\Lambda_\mathfrak q/\mathfrak q^n\Lambda_\mathfrak q
\supset
\mathfrak q\Lambda_\mathfrak q/\mathfrak q^n\Lambda_\mathfrak q
$$
\item
\label{smoothing-item-eighth-resolve}
借助给定同态
$k[y_1, \ldots, y_m]_\mathfrak p/\mathfrak p^n k[y_1, \ldots, y_m]_\mathfrak p
\to D'$ 以及赋值 $t_i\mapsto t_i$，赋予环
$D'[t_1,\ldots,t_d]$ 一个
$R_\mathfrak p/\mathfrak p^nR_\mathfrak p$-代数结构。只需找到分解
$$
B \otimes_{k[x_1, \ldots, x_d]} (R/\mathfrak p^nR)_\mathfrak p
\to D'[t_1, \ldots, t_d] \to
\Lambda_\mathfrak q/\mathfrak q^n\Lambda_\mathfrak q
$$
，其中第二个箭头把 $t_i$ 映到 $\delta_i$，并诱导给定同态
$D'\to\Lambda_\mathfrak q/\mathfrak q^n\Lambda_\mathfrak q$。
\item
\label{smoothing-item-done}
由上面对 $D'$ 的选择，这样的分解存在。
\end{enumerate}
下面逐步给出论证；仅引入记号的步骤不再说明。

\medskip\noindent
关于 (\ref{smoothing-item-power})。这是可行的，因为 $\mathfrak q$ 是包含
$\mathfrak h_A=\sqrt{H_{A/k}\Lambda}$ 的极小素理想。

\medskip\noindent
关于 (\ref{smoothing-item-strictly-standard})。注意，$A_{a_i}$ 在 $k$ 上光滑。
因此，$B_{a_j}$ 同构于 $A_{a_j}[x_1,\ldots,x_d]$ 上的多项式代数，
从而在 $k[x_1,\ldots,x_d]$ 上光滑。于是 $B_{x_i}$ 在
$k[x_1,\ldots,x_d]$ 上光滑。由引理
\ref{smoothing-lemma-improve-presentation}，$C_{x_i}$ 在
$k[x_1,\ldots,x_d]$ 上光滑，且其微分模有限自由。因此，由引理
\ref{smoothing-lemma-compare-standard}，$x_i$ 的某个幂在 $C$ 中相对于
$k[x_1,\ldots,x_n]$ 是严格标准元。

\medskip\noindent
关于 (\ref{smoothing-item-NP})。这由应用引理 \ref{smoothing-lemma-solution-modulo} 得出。

\medskip\noindent
关于 (\ref{smoothing-item-choose-pii})。由于
$k[y_1, \ldots, y_m]_\mathfrak p \to \Lambda_\mathfrak q$
依构造是平坦的，且
$\mathfrak p\Lambda_\mathfrak q=\mathfrak q\Lambda_\mathfrak q$，
由《代数》引理 \ref{algebra-lemma-dimension-base-fibre-equals-total}，
有 $\dim(k[y_1,\ldots,y_m]_\mathfrak p)=d$。因而可以找到
$\pi_1,\ldots,\pi_d\in\Lambda$，它们在 $\Lambda_\mathfrak q$
中映成一组正则参数系。

\medskip\noindent
关于 (\ref{smoothing-item-modify-pii})。由《代数》引理
\ref{algebra-lemma-regular-ring-CM}，序列
$\pi_1,\ldots,\pi_d$ 的任意排列都是
$k[y_1,\ldots,y_m]_\mathfrak p$ 中的正则序列。因此
$\gamma_1=\pi_1t_1,\ldots,\gamma_d=\pi_dt_d$ 是
$R_\mathfrak p=k[y_1,\ldots,y_m]_\mathfrak p[t_1,\ldots,t_d]$
中的正则序列，见《代数》引理
\ref{algebra-lemma-regular-sequence-in-polynomial-ring}。
令 $S=k[y_1,\ldots,y_m]\setminus\mathfrak p$，于是
$R_\mathfrak p=S^{-1}R$。注意，即使把各个 $\pi_i$ 乘以 $S$ 中元素，
$\pi_1,\ldots,\pi_d$ 和 $\gamma_1,\ldots,\gamma_d$ 仍是正则序列。
假设
$$
\text{Ann}_{R/(\gamma_1^e, \ldots, \gamma_{i - 1}^e)R}(\gamma_i)
=
\text{Ann}_{R/(\gamma_1^e, \ldots, \gamma_{i - 1}^e)R}(\gamma_i^2)
$$
对某个 $t\in\{0,\ldots,d\}$ 及 $i=1,\ldots,t$ 成立。由《代数》
引理 \ref{algebra-lemma-regular-sequence-powers}，
$\gamma_1^e,\ldots,\gamma_t^e,\gamma_{t+1}$ 是 $S^{-1}R$
中的正则序列，故
$$
\text{Ann}_{S^{-1}R/(\gamma_1^e, \ldots, \gamma_{i - 1}^e)}(\gamma_i) =
\text{Ann}_{S^{-1}R/(\gamma_1^e, \ldots, \gamma_{i - 1}^e)}(\gamma_i^2).
$$
因此，
$$
\text{Ann}_{R/(\gamma_1^e, \ldots, \gamma_t^e)R}(\gamma_{t + 1})
=
\text{Ann}_{R/(\gamma_1^e, \ldots, \gamma_t^e)R}(\gamma_{t + 1}^2)
$$
由引理 \ref{smoothing-lemma-ogoma}，取某个 $s\in S$ 并以
$s\pi_{t+1}$ 替换 $\pi_{t+1}$ 后，上式成立。对 $t$ 归纳，
即得到所需序列。

\medskip\noindent
关于 (\ref{smoothing-item-choose-deltai})。令 $S=\Lambda\setminus\mathfrak q$，
于是 $\Lambda_\mathfrak q=S^{-1}\Lambda$；并置
$\bar\Lambda=\Lambda_\mathfrak q/\mathfrak q^n\Lambda_\mathfrak q$。
假设已有 $t\in\{0,\ldots,d\}$、$\delta_1,\ldots,\delta_t\in S$，
以及 (\ref{smoothing-item-choose-deltai}) 中的分解
$D\to D'\to\bar\Lambda$，使 (a)、(b)、(c) 对
$i=1,\ldots,t$ 成立。由 (\ref{smoothing-item-power})，
$\mathfrak q^N\Lambda_\mathfrak q\subset H_{A/k}\Lambda_\mathfrak q$，
所以 $\pi_{t+1}^N\in H_{A/k}\Lambda_\mathfrak q$，进而
$\pi_{t+1}^N\in H_{A/k}\bar\Lambda$。由于
$D'\to\bar\Lambda$ 忠实平坦，由《代数》引理
\ref{algebra-lemma-faithfully-flat-universally-injective} 又有
$\pi_{t+1}^N\in H_{A/k}D'$。回忆 $H_{A/k}=(a_1,\ldots,a_t)$。
设在 $D'$ 中 $\pi_{t+1}^N=\sum a_jd_j$，并选取
$c_j\in\Lambda_\mathfrak q$ 提升 $d_j\in D'$。则
$\pi_{t+1}^N=\sum c_ja_j+\epsilon$，其中
$\epsilon\in\mathfrak q^n\Lambda_\mathfrak q\subset
\mathfrak q^{n-N}H_{A/k}\Lambda_\mathfrak q$。对某些
$c'_j\in\mathfrak q^{n-N}\Lambda_\mathfrak q$，写
$\epsilon=\sum a_jc'_j$。于是
$\pi_{t+1}^{2N}=\sum(\pi_{t+1}^Nc_j+\pi_{t+1}^Nc'_j)a_j$。
注意，$\pi_{t+1}^Nc'_j$ 在 $\bar\Lambda$ 中映到零；这个虽显然却
关键的观察将确保稍后 (a) 成立。
现在选取 $s\in S$，使存在 $\mu_{t + 1j}\in\Lambda$，一方面在
$S^{-1}\Lambda$ 中有
$\pi_{t+1}^Nc_j+\pi_{t+1}^Nc'_j=\mu_{t + 1j}/s^{2N}$，
另一方面在 $\Lambda$ 中有
$(s\pi_{t+1})^{2N}=\sum\mu_{t + 1j}a_j$（略去一个小细节）。
还可把 $s$ 替换为它的某个幂并扩大 $D'$，使 $s$ 的像属于 $D'$。
在这些选择下，$\mu_{t + 1j}$ 映到 $s^{2N}d_j$，后者属于 $D'$。
注意，由我们对 $\varphi$ 的选择，$\pi_1,\ldots,\pi_d$ 是
$S^{-1}\Lambda$ 中的一组正则参数序列。因此，由《代数》引理
\ref{algebra-lemma-regular-ring-CM}，它们构成
$\Lambda_\mathfrak q$ 中的正则序列。再由《代数》引理
\ref{algebra-lemma-regular-sequence-powers}，
${\pi'}_1^e,\ldots,{\pi'}_t^e,s\pi_{t+1}$ 是
$S^{-1}\Lambda$ 中的正则序列。所以
$$
\text{Ann}_{S^{-1}\Lambda/({\pi'}_1^e, \ldots, {\pi'}_t^e)}(s\pi_{t + 1}) =
\text{Ann}_{S^{-1}\Lambda/({\pi'}_1^e, \ldots, {\pi'}_t^e)}((s\pi_{t + 1})^2).
$$
因此，可以应用引理 \ref{smoothing-lemma-ogoma} 找到 $s'\in S$，使得
$$
\text{Ann}_{\Lambda/({\pi'}_1^e, \ldots, {\pi'}_t^e)}((s')^qs\pi_{t + 1})
=
\text{Ann}_{\Lambda/({\pi'}_1^e, \ldots, {\pi'}_t^e)}(((s')^qs\pi_{t + 1})^2).
$$
对任意 $q>0$ 都成立。由引理 \ref{smoothing-lemma-enlarge-solution-modulo}，
可以选择 $q$ 并扩大 $D'$，使 $(s')^q$ 的像属于 $D'$。
置 $\delta_{t+1}=(s')^qs$，便得 (a)、(b)、(c) 对
$i=1,\ldots,t+1$ 成立。对于 (a)，注意可取
$\lambda_{t + 1j}=(s')^{2Nq}\mu_{t + 1j}$。对 $t$ 归纳即得结论。

\medskip\noindent
关于 (\ref{smoothing-item-first-solve})。依构造，
$H_{(C\otimes_{k[x_1,\ldots,x_d]}R)/R}\Lambda$ 的根包含
$\mathfrak h_A$。事实上，元素 $a_j\in H_{A/k}$ 映到
$H_{B/k[x_1,\ldots,x_n]}$ 的元素，继而映到
$H_{C/k[x_1,\ldots,x_n]}$ 的元素，所以 $a_j\otimes1$ 映到
$H_{C\otimes_{k[x_1,\ldots,x_d]}R/R}$ 的元素。此外，若有
$C\otimes_{k[x_1,\ldots,x_n]}R\to T\to\Lambda$，它是下列情形的解：
$$
R
\to
C \otimes_{k[x_1, \ldots, x_d]} R
\to
\Lambda \supset \mathfrak q
$$
，则因 $R$ 在 $k$ 上光滑，有 $H_{T/R}\subset H_{T/k}$。
所以 $T$ 也是原情形 $k\to A\to\Lambda\supset\mathfrak q$ 的解。

\medskip\noindent
关于 (\ref{smoothing-item-second-resolve})。把引理 \ref{smoothing-lemma-lift-solution}
应用于
$R \to C \otimes_{k[x_1, \ldots, x_d]} R
\to \Lambda \supset \mathfrak q$ 以及元素序列
$\gamma_1^c,\ldots,\gamma_d^c$ 即得。注意，由于 $x_i^c$
在 $C$ 中相对于 $k[x_1,\ldots,x_d]$ 是严格标准元，引理
\ref{smoothing-lemma-strictly-standard-base-change} 表明 $\gamma_i^c$ 在
$C\otimes_{k[x_1,\ldots,x_d]}R$ 中相对于 $R$ 是严格标准元。
引理 \ref{smoothing-lemma-lift-solution} 的另一个假设由步骤
(\ref{smoothing-item-modify-pii}) 和 (\ref{smoothing-item-choose-deltai}) 成立。

\medskip\noindent
关于 (\ref{smoothing-item-third-resolve})。把引理
\ref{smoothing-lemma-delocalize-height-zero} 应用于
(\ref{smoothing-item-second-resolve}) 中的情形。以下论证中的目标环均为
Artin 局部环，因此我们所求的是经由源环上某个光滑代数 $T$ 的分解。

\medskip\noindent
关于 (\ref{smoothing-item-fifth-resolve})。假设
$C\otimes_{k[x_1,\ldots,x_d]}(R/JR)_\mathfrak p\to
T\to\Lambda_\mathfrak q/J\Lambda_\mathfrak q$ 是下列情形的解：
$$
(R/JR)_\mathfrak p \to
C \otimes_{k[x_1, \ldots, x_d]} (R/JR)_\mathfrak p \to
\Lambda_\mathfrak q/J\Lambda_\mathfrak q
\supset
\mathfrak q\Lambda_\mathfrak q/J\Lambda_\mathfrak q
$$
则
$C\otimes_{k[x_1,\ldots,x_d]}(R/I)_\mathfrak r\to T_\mathfrak r\to
\Lambda_\mathfrak q/I\Lambda_\mathfrak q$
是 (\ref{smoothing-item-third-resolve}) 中情形的解。

\medskip\noindent
关于 (\ref{smoothing-item-sixth-resolve})。$n=N+dc$ 足够大，因而
$\mathfrak p^nk[y_1,\ldots,y_m]_\mathfrak p\subset J_\mathfrak p$
且 $\mathfrak q^n\Lambda_\mathfrak q\subset J\Lambda_\mathfrak q$。
因此，若有
$C \otimes_{k[x_1, \ldots, x_d]} (R/\mathfrak p^nR)_\mathfrak p \to
T \to \Lambda_\mathfrak q/\mathfrak q^n\Lambda_\mathfrak q$
作为 (\ref{smoothing-item-fifth-resolve} 的一个解，则可取 $T/JT$ 作为
(\ref{smoothing-item-sixth-resolve}) 的解。

\medskip\noindent
关于 (\ref{smoothing-item-seventh-resolve})。这是因为在 $R$-代数范畴中
已有截面 $C\to B$。

\medskip\noindent
关于 (\ref{smoothing-item-eighth-resolve})。这是因为 $D'$ 在 Artin 局部环
$k[y_1, \ldots, y_m]_\mathfrak p/\mathfrak p^n k[y_1, \ldots, y_m]_\mathfrak p$
上本质光滑，并且
$$
R_\mathfrak p/\mathfrak p^nR_\mathfrak p =
k[y_1, \ldots, y_m]_\mathfrak p/
\mathfrak p^n k[y_1, \ldots, y_m]_\mathfrak p[t_1, \ldots, t_d].
$$
因此，$D'[t_1,\ldots,t_d]$ 是光滑
$R_\mathfrak p/\mathfrak p^nR_\mathfrak p$-代数的滤过余极限，并且
$B \otimes_{k[x_1, \ldots, x_d]} (R_\mathfrak p/\mathfrak p^nR_\mathfrak p)$
经由其中一个代数分解。

\medskip\noindent
关于 (\ref{smoothing-item-done})。证明的最后一个微妙之处在于，不能直接使用
$B\to D'$ 这个把 $x_i$ 映到 $\pi_i'$ 在 $D'$ 中的像、把
$z_{ij}$ 映到 $\lambda_{ij}$ 在 $D'$ 中之像的同态，因为需要图表
$$
\xymatrix{
B \ar[r] & D'[t_1, \ldots, t_d] \\
k[x_1, \ldots, x_d] \ar[r] \ar[u] &
R_\mathfrak p/\mathfrak p^nR_\mathfrak p \ar[u]
}
$$
可交换，并且需要复合
$B \to D'[t_1, \ldots, t_d] \to
\Lambda_\mathfrak q/\mathfrak q^n\Lambda_\mathfrak q$
等于 (\ref{smoothing-item-map-B-Lambda}) 中的同态。这就要求把 $x_i$ 映到
$\pi_it_i$ 在 $D'[t_1,\ldots,t_d]$ 中的像。因此，把 $z_{ij}$ 映到
$\lambda_{ij}t_i^{2N}/\delta_i^{2N}$ 在
$D'[t_1,\ldots,t_d]$ 中的像，一切便都清楚了。
\end{proof}







\section{主定理}
\label{smoothing-section-main}

\noindent
本节汇总前面的论证。

\begin{theorem}[Popescu]
\label{smoothing-theorem-popescu}
任意诺特环之间的正则同态都是光滑环同态的滤过余极限。
\end{theorem}

\begin{proof}
由引理 \ref{smoothing-lemma-reduce-to-field}，只需对 $k\to\Lambda$ 证明结论，
其中 $\Lambda$ 是在 $k$ 上几何正则的诺特环。设
$k\to A\to\Lambda$ 是一个分解，其中 $A$ 为有限型 $k$-代数。
只需构造分解 $A\to B\to\Lambda$，其中 $B$ 为有限型且
$\mathfrak h_B=\Lambda$，见引理 \ref{smoothing-lemma-final-solve}。
因此，可以对理想 $\mathfrak h_A$ 作诺特归纳。选取
$\mathfrak q\supset\mathfrak h_A$，使 $\mathfrak q$ 是包含
$\mathfrak h_A$ 的极小素理想。现在只需消解
$k\to A\to\Lambda\supset\mathfrak q$
（其定义见情形 \ref{smoothing-situation-local} 后的正文）。若 $k$ 的特征为零，
结论由引理 \ref{smoothing-lemma-resolve-special} 得出；若 $k$ 的特征为
$p>0$，结论由引理 \ref{smoothing-lemma-resolve-general} 得出。
\end{proof}




\section{G-环的逼近性质}
\label{smoothing-section-approximation-G-rings}

\noindent
设 $R$ 为诺特局部环。此时，$R$ 是 G-环当且仅当环同态
$R\to R^\wedge$ 正则，见《代数详论》引理
\ref{more-algebra-lemma-check-G-ring-maximal-ideals}。在这种情况下，
$R$ 的 Hensel 化 $R^h$ 及严格 Hensel 化 $R^{sh}$ 都是 G-环，
见《代数详论》引理 \ref{more-algebra-lemma-henselization-G-ring}。
此外，在一个域、完备局部环、$\mathbf{Z}$ 或特征零 Dedekind 环上
本质有限型的任意代数都是 G-环，见《代数详论》命题
\ref{more-algebra-proposition-ubiquity-G-ring}。因此，有大量环可应用
下面的结果。

\medskip\noindent
设 $R$ 为环，$f_1,\ldots,f_m\in R[x_1,\ldots,x_n]$，且 $S$ 为
$R$-代数。在这种情形下，称向量 $(a_1,\ldots,a_n)\in S^n$
是{\it $S$ 中的解}，当且仅当
$$
f_j(a_1, \ldots, a_n) = 0 \text{ 于 } S \text{ 中，对 }
j = 1, \ldots, m
$$
代数几何中的一个重要问题自然是判定多项式方程组何时有解。
下面的定理说明，在诺特局部环的完备化中有解，往往足以推出
在该环的 Hensel 化中有解。

\begin{theorem}
\label{smoothing-theorem-approximation-property}
设 $R$ 为诺特局部环，且
$f_1,\ldots,f_m\in R[x_1,\ldots,x_n]$。假设
$(a_1,\ldots,a_n)\in(R^\wedge)^n$ 是 $R^\wedge$ 中的解。
若 $R$ 是 Hensel G-环，则对每个整数 $N$，都存在 $R$ 中的解
$(b_1,\ldots,b_n)\in R^n$，使
$a_i-b_i\in\mathfrak m^NR^\wedge$。
\end{theorem}

\begin{proof}
选取 $c_i\in R$，使 $a_i-c_i\in\mathfrak m^N$。选取生成元
$\mathfrak m^N=(d_1,\ldots,d_M)$，并写成
$a_i=c_i+\sum a_{i,l}d_l$。考虑多项式环 $R[x_{i,l}]$ 及元素
$$
g_j = f_j(c_1 + \sum x_{1, l} d_l , \ldots, c_n + \sum x_{n, l} d_{n, l})
\in R[x_{i, l}]
$$
方程组 $g_j=0$ 有解 $(a_{i,l})$。假设可以证明它在 $R$ 中有解
$(b_{i,l})$，则 $b_i=c_i+\sum b_{i,l}d_l$ 是 $f_j=0$ 的解，
且模 $\mathfrak m^N$ 与 $a_i$ 同余。因此，只需证明在 $R^\wedge$
上可解蕴含在 $R$ 上可解。

\medskip\noindent
令 $A\subset R^\wedge$ 为由 $a_1,\ldots,a_n$ 生成的 $R$-子代数。
我们假设了 $R$ 是 G-环，即 $R\to R^\wedge$ 正则，因此存在分解
$$
A \to B \to R^\wedge
$$
，其中 $B$ 在 $R$ 上光滑，见定理 \ref{smoothing-theorem-popescu}。
记剩余域 $\kappa=R/\mathfrak m$。它也是 $R^\wedge$ 的剩余域，
所以得到交换图
$$
\xymatrix{
B \ar[rd] \ar@{..>}[r] & R' \ar@{..>}[d] \\
R \ar[r] \ar[u] & \kappa
}
$$
由于竖直箭头光滑，《代数详论》引理
\ref{more-algebra-lemma-lift-section-smooth-morphism} 蕴含存在平展环同态
$R\to R'$，它诱导同构 $R/\mathfrak m\to R'/\mathfrak mR'$，
并存在使上图交换的 $R$-代数同态 $B\to R'$。由于 $R$ 是 Hensel 环，
$R\to R'$ 有截面，见《代数》引理
\ref{algebra-lemma-characterize-henselian}。令 $b_i\in R$ 为 $a_i$
经环同态 $A\to B\to R'\to R$ 所得的像。由于这些全是
$R$-代数同态，$(b_1,\ldots,b_n)$ 是 $R$ 中的解。
\end{proof}

\noindent
给定诺特局部环 $(R,\mathfrak m)$、平展环同态 $R\to R'$，
以及位于 $\mathfrak m$ 之上且满足
$\kappa(\mathfrak m)=\kappa(\mathfrak m')$ 的极大理想
$\mathfrak m'\subset R'$，则有包含关系
$$
R \subset R_{\mathfrak m'} \subset R^h \subset R^\wedge,
$$
，见《代数》引理 \ref{algebra-lemma-henselian-functorial-prepare}
和《代数详论》引理 \ref{more-algebra-lemma-henselization-noetherian}。

\begin{theorem}
\label{smoothing-theorem-approximation-property-variant}
设 $R$ 为诺特局部环，且
$f_1,\ldots,f_m\in R[x_1,\ldots,x_n]$。假设
$(a_1,\ldots,a_n)\in(R^\wedge)^n$ 是一个解。若 $R$ 是 G-环，
则对每个整数 $N$，都存在
\begin{enumerate}
\item 平展环同态 $R\to R'$；
\item 位于 $\mathfrak m$ 之上的极大理想 $\mathfrak m'\subset R'$；
\item $R'$ 中的解 $(b_1,\ldots,b_n)\in(R')^n$，
\end{enumerate}
使得 $\kappa(\mathfrak m)=\kappa(\mathfrak m')$ 且
$a_i-b_i\in(\mathfrak m')^NR^\wedge$。
\end{theorem}

\begin{proof}
可以由定理 \ref{smoothing-theorem-approximation-property} 推出本定理：
由《代数详论》引理 \ref{more-algebra-lemma-henselization-G-ring}，
Hensel 化 $R^h$ 是 G-环；再把 $R^h$ 写成平展扩张 $R'$ 的有向余极限即可。
不过，我们改为在当前情形下重做上一条定理的证明。

\medskip\noindent
选取 $c_i\in R$，使 $a_i-c_i\in\mathfrak m^N$。选取生成元
$\mathfrak m^N=(d_1,\ldots,d_M)$，并写成
$a_i=c_i+\sum a_{i,l}d_l$。考虑多项式环 $R[x_{i,l}]$ 及元素
$$
g_j = f_j(c_1 + \sum x_{1, l} d_l , \ldots, c_n + \sum x_{n, l} d_{n, l})
\in R[x_{i, l}]
$$
方程组 $g_j=0$ 有解 $(a_{i,l})$。假设可以证明：对某个平展环同态
$R\to R'$，存在极大理想 $\mathfrak m'$ 满足
$\kappa(\mathfrak m)=\kappa(\mathfrak m')$，且 $g_j=0$ 在 $R'$
中有解 $(b_{i,l})$。则 $b_i=c_i+\sum b_{i,l}d_l$ 是 $f_j=0$
的解，并且模 $(\mathfrak m')^N$ 与 $a_i$ 同余。因此，只需证明
$R^\wedge$ 上可解蕴含在某个平展环扩张上可解，而且该扩张在
$\mathfrak m$ 上方的某个素理想处诱导平凡剩余域扩张。

\medskip\noindent
令 $A\subset R^\wedge$ 为由 $a_1,\ldots,a_n$ 生成的 $R$-子代数。
我们假设了 $R$ 是 G-环，即 $R\to R^\wedge$ 正则，故存在分解
$$
A \to B \to R^\wedge
$$
，其中 $B$ 在 $R$ 上光滑，见定理 \ref{smoothing-theorem-popescu}。
记剩余域 $\kappa=R/\mathfrak m$。它也是 $R^\wedge$ 的剩余域，
所以得到交换图
$$
\xymatrix{
B \ar[rd] \ar@{..>}[r] & R' \ar@{..>}[d] \\
R \ar[r] \ar[u] & \kappa
}
$$
由于竖直箭头光滑，《代数详论》引理
\ref{more-algebra-lemma-lift-section-smooth-morphism} 蕴含存在平展环同态
$R\to R'$，它诱导同构 $R/\mathfrak m\to R'/\mathfrak mR'$，
并存在使上图交换的 $R$-代数同态 $B\to R'$。令 $b_i\in R'$
为 $a_i$ 经环同态 $A\to B\to R'$ 所得的像。由于这些全是
$R$-代数同态，$(b_1,\ldots,b_n)$ 是 $R'$ 中的解。
\end{proof}

\begin{example}
\label{smoothing-example-describe-henselian}
设 $(R,\mathfrak m)$ 为诺特局部环，其 Hensel 化为 $R^h$。
完备化之间的同态 $R^\wedge\to(R^h)^\wedge$ 是同构，见《代数详论》
引理 \ref{more-algebra-lemma-henselization-noetherian}。由于 $R^h$
也是诺特环（同上），可将 $R^h$ 看作其完备化的子环
（因为完备化忠实平坦）。由此，可把 $R^h$ 视为 $R^\wedge$ 的子环。

\medskip\noindent
来考察 $R^\wedge$ 中哪些元素属于 $R^h$。为简单起见，假设 $R$
为整环，其分式域为 $K$。显然，$R^h$ 的每个元素 $f$ 都在 $R$
上代数；这意味着存在形如
$a_nf^n+\ldots+a_1f+a_0=0$ 的方程，其中 $a_i\in R$、$n>0$ 且
$a_n \not = 0$。

\medskip\noindent
反之，假设 $f\in R^\wedge$、$n\in\mathbf{N}$ 且
$a_0,\ldots,a_n\in R$，其中 $a_n \not = 0$，并满足
$a_nf^n+\ldots+a_1f+a_0=0$。若 $R$ 是 G-环，则对每个 $N>0$，
都存在 $g\in R^h$，使
$a_ng^n+\ldots+a_1g+a_0=0$ 且
$f-g\in\mathfrak m^NR^\wedge$，见定理
\ref{smoothing-theorem-approximation-property-variant}。我们希望推出当
$N\gg0$ 时 $f=g$。若非如此，就会在 $R^h$ 中找到 $P(T)$ 的无穷多个
根 $g$。这是不可能的，因为 (1) $R^h\subset R^h\otimes_RK$，
且 (2) $R^h\otimes_RK$ 是 $K$ 的有限多个域扩张的乘积。事实上，
$R\to K$ 为单射而 $R\to R^h$ 平坦，故
$R^h\to R^h\otimes_RK$ 为单射；(2) 则由《代数详论》引理
\ref{more-algebra-lemma-fibres-henselization} 得出。

\medskip\noindent
结论：若 $R$ 是分式域为 $K$ 的诺特局部整环且为 G-环，
则 $R^h\subset R^\wedge$ 恰为所有在 $K$ 上代数的元素所成的集合。
\end{example}

\noindent
下面是本节主定理的另一个变体。

\begin{lemma}
\label{smoothing-lemma-approximation-property-variant}
设 $R$ 为诺特环，$\mathfrak p\subset R$ 为素理想，且
$f_1,\ldots,f_m\in R[x_1,\ldots,x_n]$。假设
$(a_1,\ldots,a_n)\in((R_\mathfrak p)^\wedge)^n$ 是一个解。
若 $R_\mathfrak p$ 是 G-环，则对每个整数 $N$，都存在
\begin{enumerate}
\item 平展环同态 $R\to R'$；
\item 位于 $\mathfrak p$ 之上的素理想 $\mathfrak p'\subset R'$；
\item $R'$ 中的解 $(b_1,\ldots,b_n)\in(R')^n$，
\end{enumerate}
使得 $\kappa(\mathfrak p)=\kappa(\mathfrak p')$ 且
$a_i-b_i\in(\mathfrak p')^N(R'_{\mathfrak p'})^\wedge$。
\end{lemma}

\begin{proof}
由定理 \ref{smoothing-theorem-approximation-property-variant}，可以在某个
$R_\mathfrak p$ 上平展的环 $R''$ 中找到解
$(b'_1,\ldots,b'_n)$；并有位于 $\mathfrak p$ 之上的素理想
$\mathfrak p''$，满足 $\kappa(\mathfrak p)=\kappa(\mathfrak p'')$ 且
$a_i-b'_i\in(\mathfrak p'')^N(R''_{\mathfrak p''})^\wedge$。
对某个平展 $R$-代数 $R'$，可以写成
$R''=R'\otimes_RR_\mathfrak p$（见《代数》引理
\ref{algebra-lemma-etale}）。必要时把 $R'$ 替换为一个主局部化，
便可假设 $(b'_1,\ldots,b'_n)$ 来自 $R'$ 中的解
$(b_1,\ldots,b_n)$。置 $\mathfrak p'=R'\cap\mathfrak p''$，则
$R''_{\mathfrak p''}=R'_{\mathfrak p'}$，证明完成。
\end{proof}



\section{Hensel 对的逼近}
\label{smoothing-section-approximation-pairs}

\noindent
可以把第 \ref{smoothing-section-approximation-G-rings} 节的讨论推广到 Hensel 对。
Hensel 对的定义见《代数详论》第
\ref{more-algebra-section-henselian-pairs} 节。

\begin{lemma}
\label{smoothing-lemma-henselian-pair}
\begin{slogan}
Hensel 对的逼近。
\end{slogan}
设 $(A,I)$ 为 Hensel 对，其中 $A$ 是诺特环。令 $A^\wedge$
为 $A$ 的 $I$-进完备化。假设下列条件至少有一个成立：
\begin{enumerate}
\item $A\to A^\wedge$ 是正则环同态；
\item $A$ 是诺特 G-环；
\item $(A,I)$ 是某个对 $(B,J)$ 的 Hensel 化
（《代数详论》，引理 \ref{more-algebra-lemma-henselization}），
其中 $B$ 是诺特 G-环。
\end{enumerate}
给定 $f_1,\ldots,f_m\in A[x_1,\ldots,x_n]$ 以及
$\hat{a}_1,\ldots,\hat{a}_n\in A^\wedge$，满足对
$j=1,\ldots,m$ 有 $f_j(\hat{a}_1,\ldots,\hat{a}_n)=0$。
则对每个 $N\geq1$，都存在 $a_1,\ldots,a_n\in A$，使
$\hat{a}_i-a_i\in I^N$，且对 $j=1,\ldots,m$ 有
$f_j(a_1,\ldots,a_n)=0$。
\end{lemma}

\begin{proof}
由《代数详论》引理
\ref{more-algebra-lemma-henselization-pair-G-ring}，(3) 蕴含 (2)；
由《代数详论》引理
\ref{more-algebra-lemma-map-G-ring-to-completion-regular}，(2) 蕴含 (1)。
因此，只需在 $A\to A^\wedge$ 为正则环同态的情形证明本引理。

\medskip\noindent
令 $\hat{a}_1,\ldots,\hat{a}_n$ 如引理陈述所述。由定理
\ref{smoothing-theorem-popescu}，可以找到分解 $A\to B\to A^\wedge$，
其中 $A\to P$ 光滑，并有 $b_1,\ldots,b_n\in B$ 满足
$f_j(b_1,\ldots,b_n)=0$ 于 $B$ 中。记 $\sigma$ 为复合
$B\to A^\wedge\to A/I^N$。由《代数详论》引理
\ref{more-algebra-lemma-lift-section-smooth-morphism}，可以找到平展环同态
$A\to A'$，它诱导同构 $A/I^N\to A'/I^NA'$，并有提升 $\sigma$ 的
$A$-代数同态 $\tilde\sigma:B\to A'$。由于 $(A,I)$ 是 Hensel 对，
存在 $A$-代数同态 $\chi:A'\to A$，见《代数详论》引理
\ref{more-algebra-lemma-characterize-henselian-pair}。
置 $a_i=\chi(\tilde\sigma(b_i))$，便得到一个解。
\end{proof}
